% ch2.tex -- Module 8, Chapter 2: Structure of the atom (electrons, nucleus, neutrons)
% Dot points: M8.IQ2.01, M8.IQ2.02 (NSW Physics Stage 6 Syllabus 2017).
% Figures: TikZ inline, plus figures/ch2_alpha_paths.pdf (computed Coulomb paths).
%
% Chapter macro (prefix chTwo): \chTwoNuc{A}{Z}{X} nuclide symbol, A above Z, both
% right-aligned before X -- the same layout as \chOneNuc (ch1) and \chFourNuc (ch4).
\newcommand\chTwoNuc[3]{\ensuremath{\mathord{%
  \setbox0=\hbox{$\scriptstyle#1$}\setbox2=\hbox{$\scriptstyle#2$}%
  \ifdim\wd0<\wd2 \setbox0=\hbox to\wd2{\hss\box0}\else\setbox2=\hbox to\wd0{\hss\box2}\fi
  {}^{\box0}_{\box2}\mathrm{#3}}}}
\chapter[Structure of the Atom]{Structure of the Atom: Electrons, the Nucleus and Neutrons}\label{ch2:chap}
\covers{M8.IQ2.01, M8.IQ2.02}

\inquiry{How is it known that atoms are made up of protons, neutrons and electrons?}

\begin{goals}
  \item explain how cathode-ray experiments and Thomson's experiment showed that
        electrons exist, and use $v = \dfrac{E}{B}$, $\dfrac{q}{m} = \dfrac{E}{B^2 r}$
        and (for a school e/m tube) $\dfrac{q}{m} = \dfrac{2V}{B^2 r^2}$
  \item explain how Millikan showed that charge is quantised, and use
        $q = \dfrac{mgd}{V}$
  \item explain how the Geiger--Marsden experiment led to Rutherford's nuclear
        model, and assess the model
  \item explain how Chadwick used conservation laws to discover the neutron
  \item assess each experiment, including its validity, reliability and accuracy
\end{goals}

\begin{syllabus}
  \item M8.IQ2.01: investigate, assess and model the evidence for the electron:
        cathode rays, Thomson's charge-to-mass experiment, Millikan's oil drop experiment
  \item M8.IQ2.02: investigate, assess and model the evidence for the nuclear
        model: Geiger--Marsden, Rutherford's model, Chadwick's neutron
\end{syllabus}

In 1890 most scientists thought an atom was a tiny solid ball that could not be
split. Over about 40 years, a chain of experiments found the three kinds of
particle that make up atoms (\cref{ch2:fig:timeline}). Each result changed the scientists'
\term{model} of the atom: a simple picture that explains what we observe and
predicts new results. When a result does not fit, a better model replaces the old one.

\begin{figure}
  \centering
  \begin{tikzpicture}[x=1cm, y=1cm,
      ev/.style={font=\footnotesize, align=center, text width=21mm, inner sep=1pt,
        text=hscink}]
    \draw[line width=1.4pt, hscbrand!50, -{Stealth[length=3mm, width=2.4mm]}]
      (0.05,0) -- (14.95,0);
    \foreach \x in {1.2,3.0,4.8,6.6,8.4,10.2,12.0,13.8}
      \fill[hscbrand] (\x,0) circle (2.6pt);
    \node[ev, above=5pt] at (1.2,0)  {\textbf{1870s--1890s}\\cathode rays: particles or waves?};
    \node[ev, below=5pt] at (3.0,0)  {\textbf{1897}\\Thomson finds the electron ($q/m$)};
    \node[ev, above=5pt] at (4.8,0)  {\textbf{1904}\\Thomson's plum pudding model};
    \node[ev, below=5pt] at (6.6,0)  {\textbf{1909}\\Geiger and Marsden: gold foil};
    \node[ev, above=5pt] at (8.4,0)  {\textbf{1909--1913}\\Millikan measures the charge $e$};
    \node[ev, below=5pt] at (10.2,0) {\textbf{1911}\\Rutherford's nuclear model};
    \node[ev, above=5pt] at (12.0,0) {\textbf{1919}\\Rutherford identifies the proton};
    \node[ev, below=5pt] at (13.8,0) {\textbf{1932}\\Chadwick discovers the neutron};
  \end{tikzpicture}
  \caption{The story of this chapter. Each discovery built on earlier ones.}\label{ch2:fig:timeline}
\end{figure}

%% ==================================================================
\section{Discovering the electron: cathode rays and Thomson}\label{ch2:sec:electron}
\covers{M8.IQ2.01}

\begin{bigidea}
An invisible beam crosses a near-empty glass tube and makes the glass glow where
it hits. Electric and magnetic fields bend it, so it is a stream of negative
particles, not a wave. In 1897 Thomson's charge-to-mass measurement showed that
they are far lighter than any atom: they are electrons.
\end{bigidea}

\begin{keywords}
  \kw{Cathode ray tube}{a sealed glass tube with almost no air and a metal
    electrode at each end.}
  \kw{Cathode, anode}{the negative electrode and the positive electrode.}
  \kw{Charge-to-mass ratio ($q/m$)}{a particle's charge divided by its mass.
    Unit: \si{C.kg^{-1}}.}
  \kw{Velocity selector}{crossed electric and magnetic fields that let only one
    speed go straight.}
\end{keywords}

A \term{cathode ray tube} is a sealed glass tube. Nearly all the air inside has
been pumped out. At each end there is a metal plate called an \term{electrode}.
The \term{cathode} is joined to the negative terminal of a high-voltage DC
supply. (DC means direct current: one terminal always stays $+$ and the other
always stays $-$.) The \term{anode} is joined to the positive terminal (memory
trick: \emph{A}node = \emph{A}dd, $+$).

\begin{sidebyside}[0.58]
  \centering
  \begin{tikzpicture}[x=1cm, y=1cm]
    \draw[line width=0.9pt, hscink, fill=hscvel!5, rounded corners=9pt] (0,0) rectangle (7.6,2.0);
    \fill[hscvel] (0.45,0.3) rectangle (0.65,1.5);
    \fill[hscforce] (6.95,0.3) rectangle (7.15,1.5);
    \node[font=\footnotesize, anchor=west, text=hscmuted] at (0.8,1.66) {glass tube, almost no air};
    \foreach \y in {0.6,0.9,1.2}
      \draw[hscvel!75, dashed, line width=0.8pt, -{Stealth[length=2mm]}] (0.75,\y) -- (6.85,\y);
    \draw[line width=0.7pt, hscink] (0.55,0.3) -- (0.55,-0.85) -- (2.45,-0.85);
    \draw[line width=0.7pt, hscink] (7.05,0.3) -- (7.05,-0.85) -- (5.15,-0.85);
    \draw[object] (2.45,-1.35) rectangle (5.15,-0.35);
    \node[font=\footnotesize, align=center] at (3.8,-0.85) {high-voltage\\DC supply};
    \node[font=\small\bfseries, text=hscvel] at (2.25,-0.6) {$-$};
    \node[font=\small\bfseries, text=hscforce] at (5.35,-0.6) {$+$};
    \draw[hscmuted, line width=0.4pt] (0.55,1.5) -- (0.55,2.3);
    \node[diagramlabel, above, inner sep=1pt] at (0.55,2.3) {cathode ($-$)};
    \draw[hscmuted, line width=0.4pt] (7.05,1.5) -- (7.05,2.3);
    \node[diagramlabel, above, inner sep=1pt] at (7.05,2.3) {anode ($+$)};
    \draw[hscmuted, line width=0.4pt] (5.0,1.2) -- (5.0,2.3);
    \node[diagramlabel, above, inner sep=1pt, text=hscvel] at (5.0,2.3) {cathode rays};
  \end{tikzpicture}
  \captionof{figure}{A simple cathode ray tube.}\label{ch2:fig:crt}
\nextside
  When the supply is switched on, an invisible ray travels from the cathode
  towards the anode (\cref{ch2:fig:crt}). It is called a \term{cathode ray}
  because it comes from the cathode. Where it hits the glass, the glass glows.
  This is \term{fluorescence}: giving out light as soon as energy hits a material.
\end{sidebyside}

The tube needs two things. First, a \textbf{high voltage}. It pulls tiny charged
particles out of the cathode and speeds them up. (Today we call these particles
electrons. This section shows how scientists found that out.) Modern tubes heat
the cathode to release them: \term{thermionic emission} means a hot metal
``boils off'' electrons.

Second, a \textbf{very low pressure}. Then few gas particles are in the way, so
the particles travel far without colliding and stopping.

\subsection*{Particles or waves?}
Cathode rays travel through an almost empty tube and make glass glow. Light does
both of these things too. So many scientists (mostly in Germany) thought cathode
rays were waves, like light. Others (mostly in Britain) thought they were tiny
particles. The refresher below shows why electric and magnetic fields can tell
the two ideas apart.

\begin{refresher}{Charges in electric and magnetic fields (Module 6)}
\begin{itemize}
  \item Between parallel plates: $E = \dfrac{V}{d}$, pointing from $+$ to $-$.
        Electric force: $\vec{F} = q\vec{E}$. A negative charge is pushed towards the $+$ plate.
  \item Magnetic force: $F = qv_\perp B = qvB\sin\theta$. Right hand: fingers
        along $B$, thumb along $v$, palm pushes on a \emph{positive} charge.
        Reverse the answer for a negative charge.
  \item The magnetic force is at \ang{90} to $v$, so it gives circular motion:
        $F_\mathrm{c} = \dfrac{mv^2}{r}$.
  \item In diagrams, $\times$ means a field pointing into the page (you see the
        tail of the arrow). A dot $\bullet$ means out of the page (you see its tip).
  \item Light and other electromagnetic (EM) waves have no charge. Fields do not
        bend them.
\end{itemize}
\end{refresher}

\Cref{ch2:fig:tubes} shows three tubes used to test the two ideas.
\Cref{ch2:tab:evidence} lists all the evidence.

\begin{figure}
  \centering
  \begin{tikzpicture}[x=1cm, y=1cm,
      tube/.style={line width=0.8pt, hscink, fill=hscvel!5, rounded corners=7pt},
      cath/.style={fill=hscvel}, lab/.style={font=\footnotesize, align=center}]
    % (a) Maltese cross
    \begin{scope}
      \draw[tube] (0,0) rectangle (4.4,2.0);
      \fill[cath] (0.3,0.45) rectangle (0.45,1.55);
      \foreach \y in {0.55,0.7,1.3,1.45}
        \draw[hscvel!70, line width=0.6pt] (0.5,\y) -- (4.32,\y);
      \foreach \y in {0.85,1.0,1.15}
        \draw[hscvel!70, line width=0.6pt] (0.5,\y) -- (2.15,\y);
      \fill[hscink!70] (2.15,0.78) rectangle (2.3,1.22);
      \draw[hscink!70, line width=1pt] (2.22,0.78) -- (2.22,0.15);
      \draw[hscseriesC, line width=2.4pt] (4.38,0.4) -- (4.38,0.78);
      \draw[hscseriesC, line width=2.4pt] (4.38,1.22) -- (4.38,1.6);
      \draw[hscink, line width=2.4pt] (4.38,0.78) -- (4.38,1.22);
      \draw[hscmuted, line width=0.4pt] (2.22,1.22) -- (1.6,2.02);
      \node[lab, anchor=south] at (1.6,2.02) {metal cross};
      \draw[hscmuted, line width=0.4pt] (4.38,1.15) -- (3.7,2.02);
      \node[lab, anchor=south] at (3.7,2.02) {sharp shadow};
      \node[lab, anchor=north] at (2.2,-0.08) {(a) Maltese cross};
    \end{scope}
    % (b) paddle wheel
    \begin{scope}[xshift=5.15cm]
      \draw[tube] (0,0) rectangle (4.4,2.0);
      \fill[cath] (0.3,0.45) rectangle (0.45,1.55);
      \fill[hscforce] (3.95,0.45) rectangle (4.1,1.55);
      \draw[hscink, line width=0.9pt] (0.6,0.42) -- (3.8,0.42);
      \draw[hscink, line width=0.6pt] (1.9,1.0) circle (0.5);
      \foreach \a in {0,60,...,300}
        \draw[hscink, line width=1.6pt] (1.9,1.0) -- ++(\a:0.55);
      \foreach \y in {1.3,1.45}
        \draw[hscvel!70, line width=0.6pt, -{Stealth[length=1.6mm]}] (0.5,\y) -- (1.45,\y);
      \draw[displacement] (2.5,0.22) -- (3.5,0.22);
      \node[lab, anchor=north] at (2.2,-0.08) {(b) Paddle wheel};
      \node[lab, anchor=south] at (2.2,2.02) {wheel rolls away};
    \end{scope}
    % (c) charged plates
    \begin{scope}[xshift=10.3cm]
      \draw[tube] (0,0) rectangle (4.4,2.0);
      \fill[cath] (0.3,0.45) rectangle (0.45,1.55);
      \fill[hscforce] (1.4,1.62) rectangle (2.9,1.74);
      \fill[hscvel] (1.4,0.26) rectangle (2.9,0.38);
      \node[font=\small\bfseries, text=hscforce] at (3.15,1.68) {$+$};
      \node[font=\small\bfseries, text=hscvel] at (3.15,0.32) {$-$};
      \draw[hscvel, line width=1pt] (0.5,0.95) -- (1.4,0.95)
        .. controls (1.9,0.95) and (2.4,1.017) .. (2.9,1.15) -- (4.3,1.52);
      \node[lab, anchor=north] at (2.2,-0.08) {(c) Charged plates};
      \node[lab, anchor=south] at (2.2,2.02) {beam bends towards $+$};
    \end{scope}
  \end{tikzpicture}
  \caption{Three tubes used to test cathode rays (side views). In each tube the
    cathode ($-$) is the blue bar on the left, and the rays travel from left to
    right.}\label{ch2:fig:tubes}
\end{figure}

\textbf{Two warnings.} First, in 1883 Hertz's charged plates did not bend the
beam, because his tube still had some gas in it. The beam made the gas charged.
The charged gas cancelled the field between the plates. With less gas, the beam
bent (Thomson, 1897).

Second, the paddle wheel turns mainly because the beam heats its blades, not
because of the beam's momentum. In the HSC, say the turning wheel ``was taken as
evidence that the rays carry momentum''.

\begin{table}
  \caption{Evidence from the early cathode-ray experiments.}\label{ch2:tab:evidence}
  \centering\small
  \begin{tabular}{@{}L{25mm}L{66mm}L{50mm}@{}}
    \toprule
    \colhead{Test} & \colhead{What was seen} & \colhead{What it suggested}\\
    \midrule
    Fluorescence & The rays make glass and some minerals glow. &
      Waves (ultraviolet light does this too).\\
    Maltese cross & A sharp cross-shaped shadow forms at the far end. &
      Straight-line travel. Fits \emph{both} models.\\
    Thin foil & The rays pass through very thin metal foil (Hertz, 1892). &
      Waves (particles seemed too big to pass). Later: the particles are tiny.\\
    Magnet & The beam bends the way negative charges would. & Negative particles.\\
    Charged plates & No bending at first (Hertz, 1883). With a better vacuum
      the beam bends towards the $+$ plate (Thomson, 1897). & Negative particles.\\
    Collecting charge & A metal cup placed in the beam becomes negatively
      charged (Perrin, 1895). & Negative particles.\\
    Paddle wheel & A light wheel in the beam rolls away from the cathode. &
      Particles with mass and momentum (as understood then).\\
    Speed & The rays move far slower than light (Thomson, 1894). & Particles.\\
    \bottomrule
  \end{tabular}
\end{table}

\textbf{The decision.} Electric and magnetic fields bend cathode rays. They do
not bend light or other EM waves, because these have no charge. So cathode rays
are streams of negatively charged particles.

\subsection*{Thomson's charge-to-mass experiment (1897)}
\Cref{ch2:fig:thomson} shows Thomson's tube. Two anodes with small holes speed
the rays up. Each is a \term{collimator}: a plate with a narrow hole that lets
only a thin, straight beam through. The beam then passes between parallel plates
(electric field) and two coils (magnetic field). These are \term{crossed
fields}: the fields are at \ang{90} to each other and to the beam.

\begin{figure}
  \centering
  \begin{tikzpicture}[x=1cm, y=1cm, lab/.style={font=\footnotesize, align=center, inner sep=1pt}]
    \draw[line width=0.9pt, hscink, fill=hscvel!5, rounded corners=6pt]
      (3.0,0.42) -- (0,0.42) -- (0,-0.42) -- (3.0,-0.42) -- (3.0,-1.3)
      -- (10.6,-1.3) -- (10.6,1.3) -- (3.0,1.3) -- cycle;
    \draw[hscseriesC, line width=2.6pt] (10.6,-1.15) -- (10.6,1.15);
    \fill[hscvel] (0.3,-0.28) rectangle (0.45,0.28);
    \foreach \x in {1.5,2.3} {
      \fill[hscmuted] (\x,0.07) rectangle (\x+0.12,0.38);
      \fill[hscmuted] (\x,-0.38) rectangle (\x+0.12,-0.07); }
    \draw[hscfield!55, dashed, line width=0.8pt] (6.2,0) circle (1.12);
    \foreach \x in {5.4,6.2,7.0} \foreach \y in {-0.55,0.55}
      \node[text=hscfield, font=\small] at (\x,\y) {$\times$};
    \fill[hscforce] (4.9,0.95) rectangle (7.5,1.08);
    \fill[hscvel] (4.9,-1.08) rectangle (7.5,-0.95);
    \draw[hscvel, line width=1.1pt] (0.5,0) -- (10.55,0);
    \draw[hscvel, line width=0.9pt, dashed] (4.9,0) .. controls (5.77,0) and (6.63,0.083) .. (7.5,0.25) -- (10.55,0.84);
    \draw[hscvel, line width=0.9pt, dashed] (4.9,0) .. controls (5.77,0) and (6.63,-0.083) .. (7.5,-0.25) -- (10.55,-0.84);
    \foreach \y in {-0.84,0,0.84} \fill[hscseriesC!80!black] (10.6,\y) circle (2.4pt);
    \node[lab, anchor=west] at (10.75,0.84) {$E$ only};
    \node[lab, anchor=west] at (10.75,0) {$E$ and $B$\\balanced};
    \node[lab, anchor=west] at (10.75,-0.84) {$B$ only};
    \node[lab, anchor=south] at (0.38,0.5) {cathode};
    \draw[hscmuted, line width=0.4pt] (1.56,0.38) -- (1.8,0.75);
    \draw[hscmuted, line width=0.4pt] (2.36,0.38) -- (2.1,0.75);
    \node[lab, anchor=south] at (1.95,0.72) {anodes\\(collimators)};
    \node[lab, anchor=south, text=hscforce] at (6.2,1.32) {$+$ plate};
    \node[lab, anchor=north, text=hscvel] at (6.2,-1.32) {$-$ plate};
    \draw[hscmuted, line width=0.4pt] (7.0,-0.75) -- (8.3,-1.55);
    \node[lab, anchor=north west, text=hscfield] at (8.0,-1.5) {coils: $B$ into the page};
    \node[lab, anchor=south] at (10.6,1.32) {fluorescent screen};
    \draw[line width=0.6pt] (0.38,-0.28) -- (0.38,-1.1) -- (0.85,-1.1);
    \draw[line width=0.6pt] (2.36,-0.38) -- (2.36,-1.1) -- (1.45,-1.1);
    \draw[line width=0.6pt] (1.56,-0.38) -- (1.56,-1.1);
    \draw[object] (0.85,-1.27) rectangle (1.45,-0.93);
    \node[lab, anchor=north] at (1.15,-1.32) {high voltage};
    \node[font=\small\bfseries, text=hscvel] at (0.7,-0.88) {$-$};
    \node[font=\small\bfseries, text=hscforce] at (1.98,-0.88) {$+$};
  \end{tikzpicture}
  \caption{Thomson's apparatus. The electric force pushes the electrons up; the
    magnetic force pushes them down.}\label{ch2:fig:thomson}
\end{figure}

The electric force pushes the electrons up, towards the $+$ plate. The magnetic
field $B$ points into the page. Use the right-hand rule for a positive charge
moving to the right: the force is upwards. Electrons are negative, so the
magnetic force on them is downwards. So the two forces are opposite. Thomson
worked in two steps (\cref{ch2:fig:selector}).
\begin{enumerate}
  \item \textbf{Find the speed.} Switch on both fields and adjust them until the
        beam goes straight. The two forces are now equal and opposite. It is
        like a tug-of-war between two equally strong teams: the rope does not
        move sideways. Only one speed goes straight, so the fields act as a
        \term{velocity selector}.
  \item \textbf{Bend the beam.} Switch off the electric field. The magnetic force
        is always at \ang{90} to the velocity, so the beam moves on a circular
        arc of radius $r$. Thomson found $r$ from how far the spot moved.
\end{enumerate}
Gravity does not matter. Suppose there are \qty{500}{V} across plates
\qty{2.0}{cm} apart.
\begin{align*}
  E &= \frac{V}{d} = \frac{\qty{500}{V}}{\qty{0.020}{m}} = \qty{2.5e4}{V.m^{-1}}
    && \text{field between the plates}\\
  F &= qE = \qty{1.602e-19}{C}\times\qty{2.5e4}{V.m^{-1}} = \qty{4.0e-15}{N}
    && \text{electric force}\\
  F_\mathrm{g} &= mg = \qty{9.109e-31}{kg}\times\qty{9.8}{m.s^{-2}} = \qty{8.9e-30}{N}
    && \text{weight}
\end{align*}
The electric force is about
$\dfrac{\num{4.0e-15}}{\num{8.9e-30}} = \num{4.5e14}$ times bigger than the weight.

\begin{figure}
  \centering
  \begin{tikzpicture}[x=1cm, y=1cm, lab/.style={font=\footnotesize, align=center, inner sep=1pt}]
    \begin{scope}
      \fill[hscforce] (0,2.0) rectangle (4.6,2.13);
      \fill[hscvel] (0,-0.13) rectangle (4.6,0);
      \node[font=\small\bfseries, text=hscforce, anchor=east] at (-0.05,2.06) {$+$};
      \node[font=\small\bfseries, text=hscvel, anchor=east] at (-0.05,-0.06) {$-$};
      \foreach \x in {0.35,4.25} \draw[field] (\x,1.95) -- (\x,0.05);
      \node[lab, text=hscfield, anchor=west] at (4.3,1.0) {$E$};
      \foreach \p in {(0.75,1.0),(1.45,1.0),(3.95,0.45),(3.95,1.55)}
        \node[text=hscfield, font=\small] at \p {$\times$};
      \fill[hscvel] (2.3,1.0) circle (4pt);
      \node[font=\scriptsize, text=white] at (2.3,1.0) {$-$};
      \draw[velocity] (2.48,1.0) -- (3.45,1.0) node[right, text=hscvel, inner sep=1pt] {$v$};
      \draw[force] (2.3,1.18) -- (2.3,1.85) node[pos=0.55, left, text=hscforce, inner sep=1pt] {$F_\mathrm{E} = qE$};
      \draw[force] (2.3,0.82) -- (2.3,0.15) node[pos=0.55, left, text=hscforce, inner sep=1pt] {$F_\mathrm{B} = qvB$};
      \node[lab, anchor=north] at (2.3,-0.25) {(a) both fields on: balanced};
    \end{scope}
    \begin{scope}[xshift=7.6cm]
      \foreach \p in {(0.25,0.2),(0.25,0.95),(0.25,1.7),(3.3,0.2),(3.3,0.95),(2.05,0.5)}
        \node[text=hscfield!70, font=\small] at \p {$\times$};
      \draw[hscvel, line width=1pt, dashed] (1.75,0.95) circle (1.0);
      \draw[guide] (1.75,0.95) -- ++(225:1.0) node[pos=0.55, above left, inner sep=1pt, text=hscink] {$r$};
      \fill[hscink] (1.75,0.95) circle (1.2pt);
      \fill[hscvel] (1.75,1.95) circle (4pt);
      \node[font=\scriptsize, text=white] at (1.75,1.95) {$-$};
      \draw[velocity] (1.92,1.95) -- (2.75,1.95) node[right, text=hscvel, inner sep=1pt] {$v$};
      \draw[force] (1.75,1.78) -- (1.75,1.25) node[pos=0.5, left, text=hscforce, inner sep=1pt] {$F$};
      \fill[hscvel] (2.75,0.95) circle (4pt);
      \node[font=\scriptsize, text=white] at (2.75,0.95) {$-$};
      \draw[velocity] (2.75,0.78) -- (2.75,-0.05) node[below, text=hscvel, inner sep=1pt] {$v$};
      \draw[force] (2.58,0.95) -- (2.05,0.95) node[pos=0.5, above, text=hscforce, inner sep=1pt] {$F$};
      \node[lab, anchor=north] at (1.75,-0.35) {(b) magnetic field only: a circle};
      \node[lab, anchor=west, text=hscfield] at (3.75,1.6) {$B$ into\\the page};
    \end{scope}
  \end{tikzpicture}
  \caption{(a) In crossed fields the two forces on the electron cancel. (b) With
    $B$ alone the force always points to the centre; an electron goes round
    clockwise when $B$ points into the page.}\label{ch2:fig:selector}
\end{figure}

\subsection*{A school version: the e/m tube}
Schools often use an \term{e/m tube} (or fine-beam tube). It is a
glass bulb with an \term{electron gun} inside. An electron gun is a hot cathode
and an anode that make a narrow beam of electrons. The bulb holds a little gas
at low pressure. The gas glows where the electrons pass, so you can see the
beam. Two coils outside the bulb make a uniform magnetic field.

\begin{sidebyside}[0.44]
  \centering
  \begin{tikzpicture}[x=1cm, y=1cm, lab/.style={font=\footnotesize, align=center, inner sep=1pt}]
    \draw[line width=0.9pt, hscink, fill=hscvel!5] (0,0) circle (1.75);
    \foreach \p in {(-1.25,0.35),(-1.25,-0.55),(1.3,0.35),(1.3,-0.55),(-0.45,-1.4),(0.45,-1.4),(-0.45,-0.5)}
      \node[text=hscfield, font=\small] at \p {$\times$};
    \draw[object] (-1.05,0.88) rectangle (-0.25,1.06);
    \draw[hscvel, line width=1.2pt] (0,-0.05) circle (1.02);
    \draw[velocity] (0.05,0.97) -- (0.75,0.97);
    \draw[guide] (0,-0.05) -- (-30:1.02) node[pos=0.55, below, inner sep=2pt, text=hscink] {$r$};
    \fill[hscink] (0,-0.05) circle (1.2pt);
    \node[lab, anchor=south] at (-0.65,1.9) {electron gun (voltage $V$)};
    \draw[hscmuted, line width=0.4pt] (-0.65,1.06) -- (-0.65,1.9);
    \node[lab, anchor=north] at (0,-1.85) {$B$ into the page (from two coils)};
  \end{tikzpicture}
  \captionof{figure}{An e/m tube seen from the front. The beam glows, so you can
    see its circle.}\label{ch2:fig:emtube}
\nextside
  The method has two steps (\cref{ch2:fig:emtube}).
  \begin{enumerate}
    \item \textbf{Speed up.} The gun's voltage $V$ accelerates each electron
          from rest. The electric field does work $W = qV$ on it, and this
          becomes kinetic energy: $qV = \tfrac12 mv^2$.
    \item \textbf{Bend.} The magnetic field is at \ang{90} to the beam. It bends
          the beam into a full circle of radius $r$. A scale inside the bulb
          measures the circle.
  \end{enumerate}
  There is no electric field to balance, so the speed comes from $V$, not from
  $\dfrac{E}{B}$.
\end{sidebyside}

\begin{formulacard}[notonsheet]{Velocity selector (crossed fields)}
  \cardformula{v = \frac{E}{B}}
  \begin{symbols}
    \sym{$v$}{speed of the particles that go straight}{\si{m.s^{-1}}}
    \sym{$E$}{electric field strength}{\si{V.m^{-1}}}
    \sym{$B$}{magnetic field strength}{\si{T}}
  \end{symbols}
  \inwords{when the beam goes straight, its speed is $E$ divided by $B$.}
  \whentouse{$E$ and $B$ are at \ang{90} to each other and to the beam, and the
    beam is not deflected.}
  \whennottouse{the beam bends, or only one field is on.}
  \rearranged{$E = vB$ \qquad $B = \dfrac{E}{v}$}
  \tensecond{$E = \qty{3.6e4}{V.m^{-1}}$, $B = \qty{1.2e-3}{T}$:
    $v = \dfrac{3.6\times10^{4}}{1.2\times10^{-3}} = \qty{3.0e7}{m.s^{-1}}$.}
  \derive{the beam goes straight, so magnetic force $=$ electric force:
    $qvB\sin\ang{90} = qE$ (both forces are on the sheet). Use $\sin\ang{90} = 1$
    and divide both sides by $qB$: $v = \dfrac{E}{B}$.}
  \trap{it is $E$ over $B$, not $B$ over $E$. The charge cancels, so any
    particle with this speed goes straight.}
\end{formulacard}

% layout QA (book build): this toolkit used to sit just before the velocity selector
% card; there it could not fit under the e/m tube figure and left a third of a page empty.
\begin{toolkit}{Powers of ten and rearranging}
\begin{itemize}
  \item Prefixes: k $=10^{3}$, c $=10^{-2}$, m $=10^{-3}$, so
        $\qty{6.8}{cm} = \qty{0.068}{m}$. Use the $\times10^{x}$ key.
  \item Multiplying powers of ten: add the powers.
        $10^{-19}\times10^{-17} = 10^{-36}$ (see \cref{ch1:sec:bigbang}).
  \item Dividing powers of ten: subtract the powers.
        $10^{4} \div 10^{-3} = 10^{4-(-3)} = 10^{7}$.
  \item Squaring or cubing: do it to the number, and multiply the power.
        $(1.5\times10^{-3})^2 = 1.5^2\times10^{-3\times2} = 2.25\times10^{-6}$.
        On a calculator, put brackets round the whole number before you press $x^2$.
  \item Ratios: anything on both the top and the bottom cancels:
        $\dfrac{\tfrac43\pi a^3}{\tfrac43\pi b^3} = \left(\dfrac{a}{b}\right)^3$.
  \item To rearrange, do the same thing to both sides. From $qvB = qE$, divide
        both sides by $qB$ to get $v = \dfrac{E}{B}$.
\end{itemize}
\end{toolkit}

\begin{formulacard}[notonsheet]{Charge-to-mass ratio (Thomson)}
  \cardformula{\frac{q}{m} = \frac{E}{B^2 r}}
  \begin{symbols}
    \sym{$q/m$}{charge-to-mass ratio}{\si{C.kg^{-1}}}
    \sym{$E$}{electric field used in Step 1}{\si{V.m^{-1}}}
    \sym{$B$}{magnetic field (the same in both steps)}{\si{T}}
    \sym{$r$}{radius of the arc when the electric field is off}{\si{m}}
  \end{symbols}
  \inwords{a particle with more charge per kilogram bends into a tighter circle.}
  \whentouse{a Thomson-type experiment: balance the fields, then bend the beam
    with the same $B$ alone.}
  \whennottouse{$B$ changed between the steps. If you already know $v$, use
    $\dfrac{q}{m} = \dfrac{v}{Br}$.}
  \rearranged{$r = \dfrac{E}{(q/m)\,B^2}$ \qquad $\dfrac{q}{m} = \dfrac{v}{Br}$}
  \tensecond{$E = \qty{3.00e4}{V.m^{-1}}$, $B = \qty{1.50e-3}{T}$, $r = \qty{0.0758}{m}$:
    $\dfrac{q}{m} = \dfrac{3.00\times10^{4}}{(1.50\times10^{-3})^2\times0.0758}
    = \qty{1.76e11}{C.kg^{-1}}$.}
  \derive{with $E$ off, the magnetic force is the centripetal force.
    \begin{align*}
      qvB &= \frac{mv^2}{r} && \text{$F = qvB$ and $F_\mathrm{c} = \dfrac{mv^2}{r}$ (on the sheet)}\\
      \frac{qvB}{mvB} &= \frac{mv^2}{r\,mvB} && \text{divide both sides by $mvB$}\\
      \frac{q}{m} &= \frac{v}{Br} && \text{cancel $v$, $B$ and $m$}\\
      \frac{q}{m} &= \frac{E/B}{Br} = \frac{E}{B^2 r}
        && \text{put in $v = \tfrac{E}{B}$; $B\times B = B^2$}
    \end{align*}}
  \trap{$r$ in metres (\qty{6.8}{cm} $=$ \qty{0.068}{m}). Square $B$, not $r$.}
\end{formulacard}

\begin{formulacard}[notonsheet]{Charge-to-mass ratio (school e/m tube)}
  \cardformula{\frac{q}{m} = \frac{2V}{B^2 r^2}}
  \begin{symbols}
    \sym{$q/m$}{charge-to-mass ratio}{\si{C.kg^{-1}}}
    \sym{$V$}{voltage that speeds up the electrons}{\si{V}}
    \sym{$B$}{magnetic field strength}{\si{T}}
    \sym{$r$}{radius of the circle the beam makes}{\si{m}}
  \end{symbols}
  \inwords{for the same voltage and field, particles with more charge per
    kilogram bend into a smaller circle.}
  \whentouse{particles start from rest, are sped up through a voltage $V$, and
    then move in a circle at \ang{90} to $B$.}
  \whennottouse{crossed fields set the speed (use $\dfrac{q}{m} = \dfrac{E}{B^2 r}$),
    or the beam is not at \ang{90} to $B$ (it makes a spiral).}
  \rearranged{$V = \dfrac{(q/m)\,B^2 r^2}{2}$ \qquad
    $r = \dfrac{1}{B}\sqrt{\dfrac{2V}{q/m}}$}
  \tensecond{$V = \qty{300}{V}$, $B = \qty{1.5e-3}{T}$, $r = \qty{0.039}{m}$:
    $\dfrac{q}{m} = \dfrac{2\times300}{(1.5\times10^{-3})^2\times0.039^2}
    = \qty{1.8e11}{C.kg^{-1}}$.}
  \derive{use energy for the speed, then the circle.
    \begin{align*}
      qV &= \tfrac12 mv^2 && \text{$W = qV$ and $K = \tfrac12 mv^2$ (on the sheet)}\\
      v^2 &= 2\,\frac{q}{m}\,V && \text{multiply both sides by $\tfrac{2}{m}$}\\
      v &= \frac{q}{m}\,Br && \text{from $qvB = \dfrac{mv^2}{r}$ (card above)}\\
      v^2 &= \left(\frac{q}{m}\right)^2 B^2 r^2 && \text{square both sides}\\
      \left(\frac{q}{m}\right)^2 B^2 r^2 &= 2\,\frac{q}{m}\,V && \text{the two $v^2$ are equal}\\
      \frac{q}{m} &= \frac{2V}{B^2 r^2} && \text{divide both sides by $\dfrac{q}{m}B^2 r^2$}
    \end{align*}}
  \trap{$r$ is the radius. The scale often shows the diameter of the circle, so
    halve it. Square both $B$ and $r$.}
\end{formulacard}

\begin{workedexample}{Which way must the magnetic field point?}\label{ch2:we:selector}
In a Thomson-type tube, electrons move to the right between horizontal plates.
The top plate is positive and $E = \qty{2.7e4}{V.m^{-1}}$. A magnetic field of
\qty{1.5e-3}{T} makes the beam go straight. (a)~Calculate the speed of the
electrons. (b)~State and justify the direction of the magnetic field.
\given{$E = \qty{2.7e4}{V.m^{-1}}$, \ $B = \qty{1.5e-3}{T}$; electrons move
  right; top plate $+$; beam straight}
\find{the speed $v$ and the direction of $B$}
\formula{$v = \dfrac{E}{B}$; \ the magnetic force must be opposite to the electric force}
\substitute{$v = \dfrac{\qty{2.7e4}{V.m^{-1}}}{\qty{1.5e-3}{T}} = \qty{1.8e7}{m.s^{-1}}$.
  \ The electric force on the electrons is up, so the magnetic force must be down.
  With $B$ into the page, the palm pushes a \emph{positive} charge up; electrons
  are negative, so they are pushed down \cmark.}
\answer{(a)~\final{$v = \qty{1.8e7}{m.s^{-1}}$}. \ (b)~$B$ points \textbf{into the page}.}
\checkstep{units: $\unit{1\,T} = \unit{1\,N.s.C^{-1}.m^{-1}}$ and
  $\unit{1\,V.m^{-1}} = \unit{1\,N.C^{-1}}$, so $\unit{V.m^{-1}}\div\unit{T} = \unit{m.s^{-1}}$
  \cmark. This is about 6\% of the speed of light, a typical cathode-ray speed.}
\end{workedexample}

\begin{workedexample}{Thomson's result}\label{ch2:we:qm}
In \cref{ch2:we:selector} the electric field is switched off. The same magnetic
field bends the beam into an arc of radius \qty{6.8}{cm}. (a)~Calculate $q/m$.
(b)~A hydrogen ion has $q/m = \qty{9.58e7}{C.kg^{-1}}$. Compare and explain.
(A hydrogen ion is a hydrogen atom that has lost its electron. In 1897 it was the
lightest charged particle known.)
\given{$E = \qty{2.7e4}{V.m^{-1}}$, \ $B = \qty{1.5e-3}{T}$, \
  $r = \qty{6.8}{cm} = \qty{0.068}{m}$, \ hydrogen ion: \qty{9.58e7}{C.kg^{-1}}}
\find{$q/m$, and how many times bigger it is than the hydrogen ion's value}
\formula{$\dfrac{q}{m} = \dfrac{E}{B^2 r}$}
\substitute{\eqlines{B^2 &= (\qty{1.5e-3}{T})^2 = \qty{2.25e-6}{T^2}
    \quad\text{(square $B$ first)}\\
  B^2 r &= \num{2.25e-6}\times\qty{0.068}{m} = \num{1.53e-7}
    \quad\text{(multiply by $r$)}\\
  \frac{q}{m} &= \frac{\qty{2.7e4}{V.m^{-1}}}{\num{1.53e-7}} = \qty{1.76e11}{C.kg^{-1}}
    \quad\text{(divide $E$ by $B^2 r$)}\\
  \text{ratio} &= \frac{\num{1.76e11}}{\num{9.58e7}} = 1840
    \quad\text{(compare with hydrogen)}}}
\answer{(a)~\final{$q/m = \qty{1.8e11}{C.kg^{-1}}$}. \ (b)~About 1800 times the
  hydrogen ion's value. A bigger $q/m$ means more charge per kilogram. For the
  same size of charge, the particle has about $\tfrac{1}{1800}$ of a hydrogen
  atom's mass.}
\checkstep{the modern value for the electron is \qty{1.76e11}{C.kg^{-1}} \cmark.}
\end{workedexample}

\textbf{Why Thomson's result mattered.}
\begin{itemize}
  \item $q/m$ was about 1800 times that of the hydrogen ion, the lightest thing
        then known. Suppose the charge is the same size (Millikan later
        confirmed this). Then the particle has about $\tfrac{1}{1800}$ of the
        mass of a hydrogen atom. It is a \term{subatomic particle}, which means a
        particle smaller than an atom.
  \item $q/m$ was the same for every cathode metal and every gas, so the same
        particle is in all atoms. We call it the \term{electron}.
  \item Atoms are neutral, so they must also contain positive charge
        (\cref{ch2:sec:nucleus}).
\end{itemize}

\begin{workedexample}{A school e/m tube}\label{ch2:we:emtube}
In an e/m tube, electrons start from rest and are sped up through \qty{240}{V}.
A magnetic field of \qty{1.10e-3}{T} at \ang{90} to the beam bends it into a
circle. The scale shows that the circle's diameter is \qty{9.50}{cm}.
(a)~Calculate $q/m$ for the electrons. (b)~Calculate their speed.
\given{$V = \qty{240}{V}$, \ $B = \qty{1.10e-3}{T}$, \ diameter \qty{9.50}{cm}, so
  $r = \qty{4.75}{cm} = \qty{0.0475}{m}$ (halve the diameter)}
\find{$q/m$ and the speed $v$}
\formula{$\dfrac{q}{m} = \dfrac{2V}{B^2 r^2}$, \quad then \quad
  $qV = \tfrac12 mv^2$, so $v^2 = 2\,\dfrac{q}{m}\,V$}
\substitute{\eqlines{B^2 &= (\qty{1.10e-3}{T})^2 = \qty{1.21e-6}{T^2} \quad\text{(square $B$)}\\
  r^2 &= (\qty{0.0475}{m})^2 = \qty{2.256e-3}{m^2} \quad\text{(square $r$)}\\
  B^2 r^2 &= \num{1.21e-6}\times\num{2.256e-3} = \num{2.730e-9} \quad\text{(multiply them)}\\
  \frac{q}{m} &= \frac{2\times\qty{240}{V}}{\num{2.730e-9}} = \qty{1.758e11}{C.kg^{-1}}
    \quad\text{(divide $2V$ by $B^2r^2$)}\\
  v^2 &= 2\times\qty{1.758e11}{C.kg^{-1}}\times\qty{240}{V}
    \quad\text{(speed from the energy)}\\
  &= \qty{8.439e13}{m^2.s^{-2}}\\
  v &= \sqrt{\num{8.439e13}} = \qty{9.19e6}{m.s^{-1}} \quad\text{(square root)}}}
\answer{(a)~\final{$q/m = \qty{1.76e11}{C.kg^{-1}}$}. \
  (b)~\final{$v = \qty{9.2e6}{m.s^{-1}}$}.}
\checkstep{$q/m$ matches the electron's value, \qty{1.759e11}{C.kg^{-1}} \cmark.
  The speed is about 3\% of the speed of light, so we do not need relativity.}
\end{workedexample}

\begin{mistakes}
\begin{itemize}
  \item \xmark\ ``The Maltese cross proves the rays are particles.'' \cmark\ It
        shows straight-line travel, which light also does.
  \item \xmark\ $v = \dfrac{B}{E}$. \cmark\ $v = \dfrac{E}{B}$.
  \item \xmark\ Using the right-hand rule for an electron without reversing it.
        \cmark\ Find the force on a positive charge, then reverse it.
  \item \xmark\ Using the diameter of the beam's circle as $r$. \cmark\ Halve it.
\end{itemize}
\end{mistakes}

\begin{examtip}
The 2025 HSC (Question 31, 5 marks) asked about cathode-ray experiments.
Students had to assess how their results helped build an understanding of the
existence and properties of electrons. Full marks need:
\begin{enumerate}
  \item three or more experiments, each with its result and the property it showed
        (straight lines; bent by fields, so negative; Thomson's $q/m$, so far
        lighter than atoms and in all atoms);
  \item which results were weak (fluorescence and shadows fit both models;
        Hertz's poor vacuum);
  \item a judgement: ``The field experiments were decisive, because electric and
        magnetic fields cannot bend EM waves such as light.''
\end{enumerate}
For ``Outline Thomson's method'', give both steps with $qE = qvB$ and
$qvB = \dfrac{mv^2}{r}$.
\end{examtip}

\begin{quickcheck}
  \item What does the sharp shadow of the Maltese cross show? Why could it not
        settle the particle-or-wave debate?
  \item A student lets a little air back into a working cathode ray tube.
        Predict what happens to the beam, and explain why.
  \item Crossed fields of \qty{6.0e4}{V.m^{-1}} and \qty{2.4e-3}{T} let some
        electrons through undeflected. Calculate their speed.
  \item Positive ions pass undeflected through $E = \qty{2.00e4}{V.m^{-1}}$ and
        $B = \qty{0.100}{T}$. With $E$ off, the arc radius is \qty{2.09}{cm}.
        (a)~Find $q/m$. (b)~Each ion has charge $+e$. Find its mass and identify it.
  \item In an e/m tube, electrons start from rest and are sped up through
        \qty{180}{V}. (a)~Use $W = qV$ and $K = \tfrac12 mv^2$ to find their speed.
        (b)~A magnetic field of \qty{0.90e-3}{T} at \ang{90} to the beam bends
        them into a circle. Find its radius.
\end{quickcheck}

%% ==================================================================
\section{Millikan's oil drop experiment}\label{ch2:sec:millikan}
\covers{M8.IQ2.01}

\begin{bigidea}
Millikan held tiny charged oil drops still between two charged plates. Every
drop's charge was a whole number times \qty{1.6e-19}{C}. So charge comes in fixed
packets, and one packet is the charge of one electron.
\end{bigidea}

\begin{keywords}
  \kw{Quantised}{can only have certain separate values, with nothing in between.}
  \kw{Elementary charge ($e$)}{the smallest amount of charge found on its own,
    \qty{1.602e-19}{C}: the size of an electron's charge.}
  \kw{Terminal velocity}{the steady speed of a falling object when air resistance
    balances its weight.}
\end{keywords}

Something is \term{quantised} when it can only have certain separate values.
Eggs are quantised: you can buy 6 or 7 eggs, but never 6.4 eggs. From 1909 to
1913 Millikan tested whether charge is quantised (\cref{ch2:fig:millikan}).

An \term{atomiser} (a sprayer) makes a fine mist of oil drops. As the drops rub
against the nozzle, some of them become charged. Some drops drift down through
a small hole in the top plate. The two plates are horizontal and a distance $d$
apart. There is a voltage $V$ across them. A microscope shows one drop.

Two forces act on the drop. Its weight $F_\mathrm{g} = mg$ acts down. The
electric force $qE$ acts up (for a negative drop, when the top plate is positive).
Millikan adjusted $V$ until the drop hung still. Then the net force is zero, so
the two forces are equal.

\textbf{How was the mass found?} The drops are much too tiny to put on a balance.
With the field off, a drop soon falls at a steady speed, its \term{terminal velocity}.
A heavier drop falls faster, so this speed gives its mass. (Exam questions give
you the mass.) Oil was used because water evaporates: a shrinking drop would
change its mass during the measurement.

\begin{figure}
  \centering
  \begin{tikzpicture}[x=1cm, y=1cm, lab/.style={font=\footnotesize, align=center, inner sep=1pt}]
    \draw[hscink, line width=0.6pt, fill=hscmuted!10] (1.6,2.25) rectangle (5.4,3.6);
    \fill[hscforce!80] (0.55,3.05) circle (0.38);
    \draw[hscink, line width=0.6pt, fill=hscmuted!25] (0.9,2.98) rectangle (1.6,3.12);
    \foreach \p in {(2.1,3.0),(2.6,3.3),(3.0,2.75),(3.6,3.15),(4.2,2.85),(4.7,3.35),(3.3,2.5),(4.9,2.6)}
      \fill[hscacc!80] \p circle (1.6pt);
    \node[lab, anchor=south] at (0.55,3.45) {atomiser};
    \node[lab, anchor=south] at (3.5,3.62) {oil drops};
    \fill[hscforce] (1.4,2.08) rectangle (3.35,2.22);
    \fill[hscforce] (3.65,2.08) rectangle (5.6,2.22);
    \fill[hscvel] (1.4,0.58) rectangle (5.6,0.72);
    \node[font=\small\bfseries, text=hscforce] at (1.25,2.42) {$+$};
    \node[font=\small\bfseries, text=hscvel] at (1.25,0.42) {$-$};
    \foreach \x in {1.75,2.35,4.6,5.1} \draw[field] (\x,2.03) -- (\x,0.77);
    \node[lab, text=hscfield, anchor=east] at (5.06,0.95) {$E$};
    \fill[hscacc] (3.5,1.25) circle (2.6pt);
    \node[lab, anchor=south] at (3.5,1.45) {charged drop};
    \draw[<->, hscink, line width=0.6pt] (5.42,0.72) -- (5.42,2.08) node[pos=0.45, right, inner sep=2pt] {$d$};
    \draw[object] (6.15,1.98) rectangle (8.2,2.42);
    \node[font=\footnotesize] at (7.175,2.2) {X-ray source};
    \draw[hscseriesB, line width=0.9pt, -{Stealth[length=1.8mm]}] (6.13,2.1) -- (5.62,1.8);
    \draw[hscink, line width=0.6pt, fill=hscacc!25] (-0.25,1.2) rectangle (0.25,1.5);
    \node[lab, anchor=north] at (0.0,1.15) {light};
    \draw[hscacc!70, line width=0.6pt, -{Stealth[length=1.6mm]}] (0.3,1.35) -- (1.3,1.35);
    \draw[hscink, line width=0.6pt, fill=hscmuted!20] (6.4,1.22) rectangle (7.9,1.5);
    \draw[hscink, line width=0.6pt, fill=hscmuted!35] (6.25,1.27) rectangle (6.4,1.45);
    \node[lab, anchor=south] at (7.15,1.52) {microscope};
    \draw[line width=0.6pt] (1.4,2.15) -- (1.0,2.15) -- (1.0,-0.2) -- (2.6,-0.2);
    \draw[line width=0.6pt] (3.5,0.58) -- (3.5,0.0);
    \draw[object] (2.6,-0.45) rectangle (4.4,0.0);
    \node[font=\footnotesize] at (3.5,-0.23) {adjustable $V$};
    \begin{scope}[xshift=9.4cm, yshift=1.35cm]
      \draw[force] (0,0.12) -- (0,0.95) node[right, text=hscforce, inner sep=2pt] {$F_\mathrm{E} = qE$};
      \draw[force] (0,-0.12) -- (0,-0.95) node[right, text=hscforce, inner sep=2pt] {$F_\mathrm{g} = mg$};
      \fill[hscacc] (0,0) circle (3pt);
      \node[lab, anchor=north] at (0.35,-1.15) {forces on a\\still drop};
    \end{scope}
  \end{tikzpicture}
  \caption{Millikan's apparatus (left) and the forces on a drop held still
    (right). The field points down, so it pushes a negative drop up. X-rays can
    change the charge on the drop.}\label{ch2:fig:millikan}
\end{figure}


\begin{formulacard}[notonsheet]{Millikan's balanced drop}
  \cardformula{q = \frac{mgd}{V}}
  \begin{symbols}
    \sym{$q$}{charge on the drop (size)}{\si{C}}
    \sym{$m$}{mass of the drop}{\si{kg}}
    \sym{$g$}{gravitational acceleration, \qty{9.8}{m.s^{-2}}}{\si{m.s^{-2}}}
    \sym{$d$}{distance between the plates}{\si{m}}
    \sym{$V$}{voltage across the plates}{\si{V}}
  \end{symbols}
  \inwords{a heavier drop needs more charge to float; a bigger voltage needs less.}
  \whentouse{a charged drop is held still between horizontal parallel plates.}
  \whennottouse{the drop is moving at all. Even at a steady speed, air
    resistance also acts, so $qE \ne mg$. Also not when the plates are not
    horizontal.}
  \rearranged{$V = \dfrac{mgd}{q}$ \qquad $E = \dfrac{mg}{q}$}
  \tensecond{$m = \qty{2.00e-15}{kg}$, $d = \qty{0.0100}{m}$, $V = \qty{408}{V}$:
    $q = \dfrac{2.00\times10^{-15}\times9.8\times0.0100}{408} = \qty{4.80e-19}{C}$ ($= 3e$).}
  \derive{the drop is still, so the forces balance.
    \begin{align*}
      qE &= mg && \text{electric force up $=$ weight down}\\
      q\,\frac{V}{d} &= mg && \text{use $E = \dfrac{V}{d}$ (on the sheet)}\\
      q &= \frac{mgd}{V} && \text{multiply both sides by $\dfrac{d}{V}$}
    \end{align*}}
  \trap{$d$ in metres (\qty{1.2}{cm} $=$ \qty{0.012}{m}). Write weight as
    $F_\mathrm{g} = mg$: on the sheet $W$ means work.}
\end{formulacard}

\begin{formulacard}[notonsheet]{Quantised charge}
  \cardformula{q = ne}
  \begin{symbols}
    \sym{$q$}{charge on an object (size)}{\si{C}}
    \sym{$n$}{a whole number: electrons gained or lost}{no unit}
    \sym{$e$}{elementary charge, \qty{1.602e-19}{C}}{\si{C}}
  \end{symbols}
  \inwords{every charge found on its own is a whole number of elementary charges.}
  \whentouse{to check a measured charge, or to count electrons gained or lost.}
  \whennottouse{for quarks (\cref{ch5:sec:sm}), which carry $\tfrac13 e$ or
    $\tfrac23 e$ but are never found alone.}
  \rearranged{$n = \dfrac{q}{e}$}
  \tensecond{$q = \qty{8.01e-19}{C}$: $n = \dfrac{8.01\times10^{-19}}{1.602\times10^{-19}} = 5.00$,
    so 5 electrons.}
  \derive{this formula is not built from other formulas. It is what Millikan
    found by experiment. A drop becomes charged by gaining or losing whole
    electrons. Each electron carries charge $e$, so $n$ electrons carry
    $n\times e$: $q = ne$.}
  \trap{$n$ must be close to a whole number; $n = 2.5$ means the charge is wrong.
    The sheet prints the electron charge as $-\qty{1.602e-19}{C}$; use its size.}
\end{formulacard}

\begin{workedexample}{An oil drop held still}\label{ch2:we:drop}
An oil drop of mass \qty{3.9e-15}{kg} is held still between horizontal plates
\qty{1.2}{cm} apart. The voltage is \qty{477}{V} and the top plate is positive.
Find the charge on the drop and the number of extra electrons.
\given{$m = \qty{3.9e-15}{kg}$, \ $d = \qty{1.2}{cm} = \qty{0.012}{m}$, \
  $V = \qty{477}{V}$, \ $g = \qty{9.8}{m.s^{-2}}$, \ $e = \qty{1.602e-19}{C}$}
\find{the charge $q$ and the number of electrons $n$}
\formula{$q = \dfrac{mgd}{V}$, \quad then \quad $n = \dfrac{q}{e}$}
\substitute{\eqlines{mgd &= \qty{3.9e-15}{kg}\times\qty{9.8}{m.s^{-2}}\times\qty{0.012}{m}
  \\ &= \qty{4.59e-16}{J} \quad\text{(top of the fraction)}\\
  q &= \frac{\qty{4.59e-16}{J}}{\qty{477}{V}} = \qty{9.62e-19}{C} \quad\text{(divide by $V$)}\\
  n &= \frac{\qty{9.62e-19}{C}}{\qty{1.602e-19}{C}} = 6.00 \quad\text{(count the electrons)}}}
\answer{\final{$q = -\qty{9.6e-19}{C}$} (size \qty{9.6e-19}{C}), so
  \final{6 extra electrons}.}
\checkstep{$n$ is a whole number \cmark. The drop is pulled towards the $+$
  plate, so it is negative, which fits extra electrons \cmark.}
\end{workedexample}

\begin{workedexample}{Finding $e$ from data}\label{ch2:we:data}
An oil drop simulation gives these charges (in units of \qty{e-19}{C}):
6.38, 3.22, 9.63, 4.79, 11.19, 8.03. Estimate the elementary charge.
\given{six charges, in units of \qty{e-19}{C}}
\find{the smallest unit of charge, $e$}
\formula{$q = ne$: every charge, and every gap between two charges, is a whole
  number of $e$}
\substitute{\textbf{Step 1.} Sort the charges: 3.22, 4.79, 6.38, 8.03, 9.63, 11.19.\par
  \textbf{Step 2.} Find the gap between neighbours: 1.57, 1.59, 1.65, 1.60, 1.56.
  All the gaps are about the same, so each gap is one step of $e$. (A gap about
  twice as big would be $2e$.)\par
  \textbf{Step 3.} Average gap $= (1.57+1.59+1.65+1.60+1.56)\div5 = 7.97\div5 = 1.594$.\par
  \textbf{Step 4.} Divide each charge by 1.594: 2.02, 3.01, 4.00, 5.04, 6.04,
  7.02. Round them: $n = 2, 3, 4, 5, 6, 7$.\par
  \textbf{Step 5.} Use every drop for the best value. Add the charges
  ($43.24\times10^{-19}$~C) and the $n$ values ($2+3+4+5+6+7 = 27$):
  $e = \dfrac{43.24\times10^{-19}\ \unit{C}}{27} = \qty{1.601e-19}{C}$.
  Using all six drops lets random errors partly cancel.}
\answer{\final{$e = \qty{1.60e-19}{C}$}}
\checkstep{our unrounded value, \qty{1.601e-19}{C}, is within 0.1\% of the data
  sheet value \qty{1.602e-19}{C} \cmark. Every $n$ is close to a whole number,
  so the data fit $q = ne$.}
\end{workedexample}

\textbf{The electron's mass.} For the electron, Thomson's ratio $q/m$ is written
$\dfrac{e}{m_\mathrm{e}}$ (charge $e$, mass $m_\mathrm{e}$). Thomson measured this
ratio, and Millikan measured $e$. So
$m_\mathrm{e} = e \div \dfrac{e}{m_\mathrm{e}}
= \dfrac{\qty{1.602e-19}{C}}{\qty{1.759e11}{C.kg^{-1}}} = \qty{9.11e-31}{kg}$,
the value on the data sheet. That is about $\tfrac{1}{1836}$ of a proton's mass.

\begin{mistakes}
\begin{itemize}
  \item \xmark\ $q = \dfrac{mgV}{d}$. \cmark\ $q = \dfrac{mgd}{V}$
        (units: $\unit{N.m} = \unit{J}$, and $\unit{1\,V} = \unit{1\,J.C^{-1}}$,
        so $\unit{J}\div\unit{V} = \unit{C}$).
  \item \xmark\ Giving a charge in eV. \cmark\ eV is energy; charge is in coulombs.
  \item \xmark\ ``Millikan measured one electron directly.'' \cmark\ He measured
        oil drops; the common step between their charges is $e$.
  \item \xmark\ Taking the smallest measured charge as $e$. \cmark\ In
        \cref{ch2:we:data} the smallest charge, \qty{3.22e-19}{C}, is $2e$. Use the gaps.
\end{itemize}
\end{mistakes}

\begin{examtip}
``Explain how Millikan's results show that charge is quantised'' (3 marks).
Full marks need:
\begin{enumerate}
  \item many drops were measured, each with $q = \dfrac{mgd}{V}$;
  \item every charge was a whole number times \qty{1.6e-19}{C}, with nothing in between;
  \item drops gain or lose whole electrons, so this step is the size of the
        electron's charge.
\end{enumerate}
In data questions, show the gaps or divide by $e$. Comment on any value that
does not fit.
\end{examtip}

\begin{quickcheck}
  \item A drop of mass \qty{2.4e-15}{kg} carries 3 extra electrons between plates
        \qty{8.0}{mm} apart. What voltage holds it still?\label{ch2:qc:volt}
  \item A student reports a drop charge of \qty{4.0e-19}{C}. Explain why this
        must be wrong.
  \item The plates in \cref{ch2:qc:volt} are swapped: the top plate becomes
        negative and the bottom plate positive. Describe and explain what
        happens to the drop.
\end{quickcheck}

%% ==================================================================
\section{The nuclear atom: Geiger--Marsden and Rutherford}\label{ch2:sec:nucleus}
\covers{M8.IQ2.02}

\begin{bigidea}
Alpha particles were fired at very thin gold foil. Almost all went straight
through, but a very few bounced back. So an atom is mostly empty space, with
nearly all its mass and positive charge in a tiny centre: the nucleus.
\end{bigidea}

\begin{keywords}
  \kw{Plum pudding model}{Thomson's model: spread-out positive charge with
    electrons dotted through it.}
  \kw{Alpha particle ($\alpha$)}{a helium nucleus: 2 protons and 2 neutrons, charge $+2e$.}
  \kw{Nucleus}{the tiny, dense, positive centre of an atom.}
\end{keywords}

\textbf{Thomson's model (1904).} Atoms contain electrons, but atoms are neutral.
So each atom must also hold positive charge. In the \term{plum pudding model},
this positive charge is spread evenly through the whole atom. The electrons are
dotted through it, like blueberries in a muffin (\cref{ch2:fig:models}a). (A plum
pudding is an old English dessert with fruit spread through it.) There is no nucleus.

\begin{figure}
  \centering
  \begin{tikzpicture}[x=1cm, y=1cm, lab/.style={font=\footnotesize, align=left, inner sep=1pt}]
    \begin{scope}
      \fill[hscforce!13] (0,0) circle (1.3);
      \draw[hscforce!60, line width=0.8pt] (0,0) circle (1.3);
      \foreach \p in {(-0.6,0.55),(0.15,0.8),(0.7,0.35),(-0.85,-0.2),(-0.1,0.05),(0.55,-0.5),(-0.35,-0.75)}
        \node[text=hscforce, font=\small\bfseries] at \p {$+$};
      \foreach \p in {(-0.25,0.5),(0.45,0.0),(-0.6,-0.5),(0.2,-0.85),(0.8,0.75),(-0.95,0.25),(0.15,-0.35)} {
        \fill[hscvel] \p circle (3.2pt);
        \node[text=white, font=\tiny] at \p {$-$}; }
      \draw[hscmuted, line width=0.4pt] (0.98,0.52) -- (1.75,1.05);
      \node[lab, anchor=west] at (1.75,1.05) {positive charge spread\\through the whole atom};
      \draw[hscmuted, line width=0.4pt] (0.48,0.03) -- (1.75,-0.35);
      \node[lab, anchor=west] at (1.75,-0.35) {electrons};
      \node[font=\footnotesize, anchor=north] at (1.2,-1.4) {(a) Plum pudding model (Thomson, 1904)};
    \end{scope}
    \begin{scope}[xshift=8.2cm]
      \draw[hscmuted, dashed, line width=0.6pt] (0,0) circle (1.3);
      \fill[hscforce] (0,0) circle (1.6pt);
      \foreach \p in {(-0.7,0.6),(0.55,0.9),(0.3,-0.95),(-0.95,-0.45),(0.95,-0.25),(-0.25,-0.8)} {
        \fill[hscvel] \p circle (3.2pt);
        \node[text=white, font=\tiny] at \p {$-$}; }
      \draw[hscmuted, line width=0.4pt] (0.05,0.03) -- (1.75,0.95);
      \node[lab, anchor=west] at (1.75,0.95) {tiny nucleus ($+$)};
      \node[lab, anchor=west] at (1.75,0.2) {mostly empty space};
      \draw[hscmuted, line width=0.4pt] (0.3,-0.95) -- (1.75,-0.55);
      \node[lab, anchor=west] at (1.75,-0.55) {electrons around\\the nucleus};
      \node[font=\footnotesize, anchor=north] at (1.2,-1.4) {(b) Nuclear model (Rutherford, 1911)};
    \end{scope}
  \end{tikzpicture}
  \caption{Two models of the atom. In (b) the nucleus is drawn far too big: it is
    really about \num{10000} to \num{100000} times smaller across than the atom.}\label{ch2:fig:models}
\end{figure}

\textbf{The test (1909).} Geiger and Marsden worked for Rutherford. They fired
\term{alpha particles} at gold foil about a thousand atoms thick
(\cref{ch2:fig:goldfoil}). An alpha particle is a helium nucleus with charge
$+2e$. Radioactive atoms (atoms that break apart by themselves) give them out.
A zinc sulfide screen gives a tiny flash, a \term{scintillation}, where an alpha
hits it. In a dark room, they used a microscope to count the flashes at each angle.

\begin{figure}
  \centering
  \begin{tikzpicture}[x=1cm, y=1cm, lab/.style={font=\footnotesize, align=left, inner sep=1pt}]
    \draw[hscmuted, dashed, line width=0.6pt] (0,0) circle (2.55);
    \draw[hscseriesC, line width=2.2pt] (-168:2.05) arc (-168:168:2.05);
    \fill[hscmuted!45] (-4.7,-0.5) rectangle (-3.5,0.5);
    \fill[white] (-4.0,-0.07) rectangle (-3.5,0.07);
    \fill[hscacc] (-4.15,0) circle (3pt);
    \draw[hscacc, line width=1.1pt, -{Stealth[length=2mm]}] (-3.45,0) -- (-0.08,0);
    \draw[draw=yellow!60!black, line width=2.4pt] (0,-0.55) -- (0,0.55);
    \foreach \a in {-3,0,3} \draw[hscacc, line width=0.9pt, -{Stealth[length=1.8mm]}] (0,0) -- (\a:1.95);
    \draw[hscacc, line width=0.9pt, -{Stealth[length=1.8mm]}] (0,0) -- (25:1.95);
    \draw[hscacc, line width=0.9pt, -{Stealth[length=1.8mm]}] (0,0) -- (-30:1.95);
    \draw[hscacc, line width=0.9pt, -{Stealth[length=1.8mm]}] (0,0) -- (140:1.95);
    \begin{scope}[rotate=55]
      \draw[hscink, line width=0.6pt, fill=hscmuted!20] (2.12,-0.14) rectangle (3.1,0.14);
    \end{scope}
    \node[lab, anchor=north] at (-4.1,-0.55) {alpha source\\in a lead block};
    \draw[hscmuted, line width=0.4pt] (-0.04,-0.52) -- (-0.4,-0.74);
    \node[lab, anchor=north] at (-0.6,-0.74) {thin gold foil};
    \node[lab, anchor=west] at (2.7,-0.05) {most: straight through};
    \node[lab, anchor=west] at (2.4,0.95) {some: small angle};
    \draw[hscmuted, line width=0.4pt] (140:1.75) -- (-2.25,2.0);
    \node[lab, anchor=south east] at (-2.2,1.95) {about 1 in 8000:\\more than \ang{90}};
    \node[lab, anchor=west] at (2.2,2.3) {microscope (moves\\around the foil)};
    \draw[hscmuted, line width=0.4pt] (-45:2.05) -- (2.6,-2.0);
    \node[lab, anchor=west] at (2.6,-2.0) {zinc sulfide screen};
    \node[lab, anchor=north] at (-1.9,-2.2) {vacuum chamber};
    \draw[hscmuted, line width=0.4pt] (-1.9,-2.2) -- (-130:2.55);
  \end{tikzpicture}
  \caption{The Geiger--Marsden experiment, seen from above.}\label{ch2:fig:goldfoil}
\end{figure}

\textbf{The prediction.} In the plum pudding model the positive charge is spread
thinly, so its force is weak everywhere. Electrons are about 7300 times lighter
than an alpha, so they cannot turn it (ping-pong balls cannot turn a bowling
ball). So every alpha should go through with only a tiny deflection.
\Cref{ch2:tab:goldfoil} shows what really happened.

\begin{table}
  \caption{The results, and how Rutherford explained them (1911).}\label{ch2:tab:goldfoil}
  \centering\small
  \begin{tabular}{@{}L{56mm}L{84mm}@{}}
    \toprule
    \colhead{Observation} & \colhead{Rutherford's explanation}\\
    \midrule
    Most alphas went straight through. & The atom is mostly empty space.\\
    Some were deflected by small angles. & They passed fairly close to a
      positive nucleus and were pushed sideways.\\
    A very small fraction (about 1 in 8000 in the 1909 measurements) were
      deflected by more than \ang{90}. &
      They headed almost straight at a tiny, positive nucleus holding nearly all
      the mass. The huge repulsion turned them back. This is rare because the
      nucleus is so small.\\
    \bottomrule
  \end{tabular}
\end{table}

Rutherford was amazed. He said it was like firing a heavy shell from a big gun
at tissue paper and seeing it bounce back. In 1911 he
proposed the \term{nuclear model}:
\begin{itemize}
  \item a tiny nucleus holds all the positive charge and more than 99.9\% of the mass;
  \item electrons move around the nucleus, far away from it;
  \item most of the atom is empty space.
\end{itemize}
The alphas never touch the nucleus. The electric force repels them as they come
close (\cref{ch2:fig:paths}).

\begin{figure}
  \centering
  \includegraphics{ch2_alpha_paths.pdf}
  \caption{Paths of alpha particles near one nucleus, calculated with Coulomb's
    law. The closer an alpha comes, the more it is turned.}\label{ch2:fig:paths}
\end{figure}

\textbf{Assessing Rutherford's model.} To assess a model, weigh its strengths
against its limitations, then make a judgement.
\begin{itemize}
  \item \textbf{Strengths.} It explains all three gold-foil results and why atoms
        are neutral. In 1913 Geiger and Marsden counted alphas at many angles.
        The counts matched Rutherford's predictions from Coulomb's law. They also
        changed in the predicted way with foil thickness, with the metal (its
        nuclear charge) and with the alpha speed. This was strong evidence for a
        tiny, charged nucleus.
  \item \textbf{Limitations.} (1)~An orbiting electron is accelerating. Classical
        physics says it should give out light, lose energy and spiral into the
        nucleus. (2)~It cannot explain line spectra (\cref{ch3:sec:bohr}). (3)~It does
        not say what the nucleus is made of. It also cannot explain why nuclei
        have about twice the mass of their protons (\cref{ch2:sec:neutron}).
  \item \textbf{Judgement.} The model is strongly supported for the nucleus, but
        it is incomplete for the electrons.
\end{itemize}

\textbf{Modelling it in class.} You cannot see a nucleus, but you can model the
idea. Hide an object under a board. Roll marbles under the board and record
where they come out. From the paths, work out the hidden object's shape and
size. This is indirect evidence, as in the Geiger--Marsden experiment. A
computer simulation of alpha scattering shows the same idea.

\begin{refresher}{Coulomb's law (Year 11)}
Coulomb's law gives the force between two small charges. It uses
$\varepsilon_0$, the \emph{permittivity of free space}, a constant on the data
sheet ($\num{8.854e-12}$ in SI units). Work out this part of the formula once:
\[ \frac{1}{4\pi\varepsilon_0} = \frac{1}{4\pi\times8.854\times10^{-12}}
= \qty{8.99e9}{N.m^2.C^{-2}} \]
Use this number in every Coulomb's law calculation.
\end{refresher}

\begin{formulacard}{Coulomb's law}
  \cardformula{F = \frac{1}{4\pi\varepsilon_0}\frac{q_1 q_2}{r^2}}
  \begin{symbols}
    \sym{$F$}{electric force between two charges}{\si{N}}
    \sym{$q_1$, $q_2$}{the two charges}{\si{C}}
    \sym{$r$}{distance between their centres}{\si{m}}
    \sym{$\varepsilon_0$}{permittivity of free space, \num{8.854e-12}; so
      $\frac{1}{4\pi\varepsilon_0} = \qty{8.99e9}{N.m^2.C^{-2}}$}{\si{C^2.N^{-1}.m^{-2}}}
  \end{symbols}
  \inwords{like charges push apart and unlike charges pull together; double the
    distance and the force is 4 times weaker.}
  \whentouse{two small charged particles, such as an alpha particle and a nucleus.}
  \whennottouse{between parallel plates (use $F = qE$), or for a particle inside
    a spread-out charge.}
  \rearranged{$r = \sqrt{\dfrac{q_1 q_2}{4\pi\varepsilon_0 F}}$ \qquad
    $q_2 = \dfrac{4\pi\varepsilon_0 F r^2}{q_1}$}
  \tensecond{two protons \qty{1.0e-15}{m} apart:
    $F = \dfrac{8.99\times10^{9}\times(1.602\times10^{-19})^2}{(1.0\times10^{-15})^2}
    = \qty{2.3e2}{N}$.}
  \trap{$r$ is measured centre to centre, in metres, and it is squared. An
    alpha's charge is $2e = \qty{3.204e-19}{C}$, not 2.}
\end{formulacard}

\begin{workedexample}{How small is the nucleus?}\label{ch2:we:size}
A gold atom has a radius of about \qty{1.44e-10}{m}; its nucleus, about
\qty{7.0e-15}{m}. How many times bigger is the atom's radius, and what fraction of
the atom's volume does the nucleus fill?
\given{$r_\mathrm{atom} = \qty{1.44e-10}{m}$, \ $r_\mathrm{nucleus} = \qty{7.0e-15}{m}$}
\find{$\dfrac{r_\mathrm{atom}}{r_\mathrm{nucleus}}$ and the volume fraction}
\formula{sphere volume $= \tfrac43\pi r^3$, so
  $\dfrac{V_\mathrm{nucleus}}{V_\mathrm{atom}}
  = \dfrac{\tfrac43\pi r_\mathrm{nucleus}^3}{\tfrac43\pi r_\mathrm{atom}^3}
  = \left(\dfrac{r_\mathrm{nucleus}}{r_\mathrm{atom}}\right)^3$
  ($\tfrac43\pi$ is on the top and the bottom, so it cancels)}
\substitute{\eqlines{\frac{r_\mathrm{atom}}{r_\mathrm{nucleus}} &=
  \frac{\qty{1.44e-10}{m}}{\qty{7.0e-15}{m}} = \num{2.06e4} \quad\text{(how many times wider)}\\
  \frac{r_\mathrm{nucleus}}{r_\mathrm{atom}} &=
  \frac{\qty{7.0e-15}{m}}{\qty{1.44e-10}{m}} = \num{4.86e-5} \quad\text{(flip the ratio)}\\
  \left(\frac{r_\mathrm{nucleus}}{r_\mathrm{atom}}\right)^3 &=
  4.86^3\times10^{-5\times3} \quad\text{(cube both parts)}\\
  &= 114.8\times10^{-15} = \num{1.1e-13}}}
\answer{the atom is about \final{\num{21000} times} wider; the nucleus fills
  about \final{$10^{-13}$} of its volume.}
\checkstep{a marble \qty{1}{cm} across as the nucleus gives an atom about
  \qty{200}{m} across: a marble in the middle of a stadium.}
\end{workedexample}

\begin{workedexample}{Why must the positive charge be concentrated?}\label{ch2:we:coulomb}
An alpha particle ($+2e$) moves towards a gold atom ($+79e$). Ignore the
electrons. Find the electric force on the alpha in each model.
(a)~Nuclear model: the alpha is \qty{3.0e-14}{m} from the nucleus.
(b)~Plum pudding model: the alpha is at the atom's surface
($r = \qty{1.44e-10}{m}$). The spread-out charge pushes hardest here.
\given{$q_1 = 2\times\qty{1.602e-19}{C} = \qty{3.204e-19}{C}$, \
  $q_2 = 79\times\qty{1.602e-19}{C} = \qty{1.266e-17}{C}$, \
  $\dfrac{1}{4\pi\varepsilon_0} = \qty{8.99e9}{N.m^2.C^{-2}}$}
\find{the force in each model, and the ratio}
\formula{$F = \dfrac{1}{4\pi\varepsilon_0}\dfrac{q_1 q_2}{r^2}$}
\substitute{\eqlines{q_1 q_2 &= (\qty{3.204e-19}{C})(\qty{1.266e-17}{C})
    \quad\text{(multiply the charges)}\\
  &= \qty{4.055e-36}{C^2}\\
  r_\mathrm{(a)}^2 &= (\qty{3.0e-14}{m})^2 = \qty{9.0e-28}{m^2} \quad\text{(square $r$)}\\
  r_\mathrm{(b)}^2 &= (\qty{1.44e-10}{m})^2 = \qty{2.07e-20}{m^2} \quad\text{(square $r$)}\\
  F_\mathrm{(a)} &= \frac{\qty{8.99e9}{N.m^2.C^{-2}}\times\qty{4.055e-36}{C^2}}{\qty{9.0e-28}{m^2}} = \qty{40.5}{N}\\
  F_\mathrm{(b)} &= \frac{\qty{8.99e9}{N.m^2.C^{-2}}\times\qty{4.055e-36}{C^2}}{\qty{2.07e-20}{m^2}}
  = \qty{1.76e-6}{N}}}
\answer{\final{$F_\mathrm{(a)} = \qty{41}{N}$}, \ \final{$F_\mathrm{(b)} = \qty{1.8e-6}{N}$}:
  the concentrated nucleus pushes about \num{2.3e7} times harder. Only a tiny
  nucleus can turn a fast alpha back.}
\checkstep{$F \propto \dfrac{1}{r^2}$, and
  $\left(\dfrac{1.44\times10^{-10}}{3.0\times10^{-14}}\right)^2 = \num{2.3e7}$ \cmark.}
\end{workedexample}

\begin{mistakes}
\begin{itemize}
  \item \xmark\ ``The plum pudding model has a nucleus.'' \cmark\ It has none:
        the positive charge fills the whole atom.
  \item \xmark\ ``The alphas bounced off by hitting the nucleus.'' \cmark\ The
        electric force repelled them before they touched it.
  \item \xmark\ ``Most alphas bounced back.'' \cmark\ Only about 1 in 8000 were
        turned by more than \ang{90}.
\end{itemize}
\end{mistakes}

\begin{examtip}
``Describe the Geiger--Marsden experiment and explain how it led to Rutherford's
model'' (5--6 marks). Full marks need:
\begin{enumerate}
  \item the method: alpha source, narrow beam, thin gold foil, zinc sulfide
        screen, vacuum, flashes counted at many angles (a labelled diagram helps);
  \item all three observations, each with its conclusion (\cref{ch2:tab:goldfoil});
  \item the contrast with the plum pudding prediction.
\end{enumerate}
For ``Assess Rutherford's model'', give two strengths, two limitations and a judgement.
\end{examtip}

\begin{quickcheck}
  \item Sketch and label Thomson's plum pudding atom. Use your sketch to predict
        what should happen to alpha particles fired at thin gold foil.
  \item Explain why most alpha particles passed straight through the foil, but a
        few bounced back.
  \item Calculate the force on an alpha particle ($+2e$) \qty{2.5e-14}{m} from a
        silver nucleus ($+47e$).
  \item Assess Rutherford's model of the atom. (4 marks)
\end{quickcheck}

%% ==================================================================
\section{Chadwick and the neutron}\label{ch2:sec:neutron}
\covers{M8.IQ2.02}

\begin{bigidea}
Most nuclei have about twice the mass of their protons. In 1932 Chadwick
used conservation of energy and momentum to explain a strange neutral
``radiation''. It was a stream of neutrons: particles with about a proton's mass
but no charge.
\end{bigidea}

\begin{keywords}
  \kw{Neutron}{a neutral particle in the nucleus with about the mass of a proton.}
  \kw{Mass number ($A$)}{the number of protons plus neutrons in a nucleus.}
  \kw{Atomic number ($Z$)}{the number of protons in a nucleus.}
  \kw{Elastic collision}{a collision in which the total kinetic energy stays the same.}
  \kw{Recoil}{to move off after being hit.}
\end{keywords}

\textbf{The missing mass.} In 1919 Rutherford fired alpha particles at nitrogen.
Hydrogen nuclei were knocked out. In 1920 he named them \term{protons}: positive
particles with charge $+e$. But nuclei were too heavy. A helium nucleus has
charge $+2e$, so it has 2 protons. Two protons should give 2 times the mass of
hydrogen, but helium has about 4 times. So something with mass but no charge
must also be in the nucleus. In the same year, Rutherford predicted a neutral particle.

\begin{refresher}{Nuclear symbols (Year 11)}
$\chTwoNuc{A}{Z}{X}$: $\chTwoNuc{9}{4}{Be}$ has 4 protons and $9-4 = 5$ neutrons.
Neutron $\chTwoNuc{1}{0}{n}$, proton $\chTwoNuc{1}{1}{p}$, alpha
$\chTwoNuc{4}{2}{He}$. In a reaction, the totals of $A$ and of $Z$ are the same
on both sides.
\end{refresher}

\textbf{The mystery radiation.} In 1930 Bothe and Becker fired alpha particles at
beryllium. A new radiation came out. It was very penetrating: it could pass
through thick sheets of lead. Electric and magnetic fields did not bend it, so it
had no charge. They thought it was gamma rays (very high-energy light).

In 1932 Irène Curie and Frédéric Joliot aimed this radiation at paraffin wax.
Wax contains a lot of hydrogen, so it is full of protons. Fast protons shot out
of the wax. Chadwick repeated this with other targets, such as nitrogen. He
measured the energies of the nuclei knocked out (\cref{ch2:fig:chadwick}). His
detector was an \term{ionisation chamber}. A charged particle passing through
its gas knocks electrons off the gas atoms, and the tiny current this makes is
measured.

\begin{figure}
  \centering
  \begin{tikzpicture}[x=1cm, y=1cm,
      bx/.style={draw=hscink, line width=0.7pt, fill=hscmuted!12, rounded corners=2pt,
        minimum width=18mm, minimum height=10mm, align=center, font=\footnotesize},
      lab/.style={font=\footnotesize, align=center, inner sep=1pt}]
    \node[bx] (po) at (1.1,0) {polonium\\(alpha source)};
    \node[bx] (be) at (4.85,0) {beryllium};
    \node[bx] (wax) at (8.75,0) {paraffin wax\\(lots of H)};
    \node[bx] (det) at (13.1,0) {ionisation chamber\\(detector)};
    \draw[hscacc, line width=1.1pt, -{Stealth[length=2mm]}] (po.east) -- (be.west)
      node[midway, above, lab] {alpha\\particles};
    \draw[hscmuted, line width=1.1pt, dashed, -{Stealth[length=2mm]}] (be.east) -- (wax.west)
      node[midway, above, lab] {neutral\\radiation};
    \draw[hscforce, line width=1.1pt, -{Stealth[length=2mm]}] (wax.east) -- (det.west)
      node[midway, above, lab] {protons};
    \begin{scope}[yshift=-1.75cm]
      \node[lab, anchor=west] at (-0.05,0.35) {neutron hits a proton (same mass):};
      \fill[hscmuted] (0.3,-0.25) circle (5pt); \node[lab, below=4pt] at (0.3,-0.25) {n};
      \draw[velocity] (0.5,-0.25) -- (1.3,-0.25);
      \fill[hscforce] (1.7,-0.25) circle (5pt); \node[lab, below=4pt] at (1.7,-0.25) {p};
      \node[font=\small] at (2.45,-0.25) {$\Rightarrow$};
      \fill[hscmuted] (3.2,-0.25) circle (5pt); \node[lab, below=4pt] at (3.2,-0.25) {stops};
      \fill[hscforce] (4.1,-0.25) circle (5pt);
      \draw[velocity] (4.3,-0.25) -- (5.1,-0.25);
      \node[lab, anchor=north west] at (-0.05,-0.75) {the proton gets almost all the energy};
    \end{scope}
    \begin{scope}[xshift=7.2cm, yshift=-1.75cm]
      \node[lab, anchor=west] at (-0.05,0.35) {neutron hits nitrogen (14 times the mass):};
      \fill[hscmuted] (0.3,-0.25) circle (5pt); \node[lab, below=4pt] at (0.3,-0.25) {n};
      \draw[velocity] (0.5,-0.25) -- (1.3,-0.25);
      \fill[hscseriesA!70] (1.85,-0.25) circle (10pt); \node[font=\scriptsize, text=white] at (1.85,-0.25) {N};
      \node[font=\small] at (2.7,-0.25) {$\Rightarrow$};
      \draw[velocity] (3.75,-0.25) -- (3.05,-0.25);
      \fill[hscmuted] (3.95,-0.25) circle (5pt); \node[lab, below=4pt] at (3.95,-0.25) {bounces back};
      \fill[hscseriesA!70] (5.5,-0.25) circle (10pt); \node[font=\scriptsize, text=white] at (5.5,-0.25) {N};
      \draw[velocity] (5.9,-0.25) -- (6.15,-0.25);
      \node[lab, anchor=north west] at (-0.05,-0.75) {nitrogen gets only about a quarter of the energy};
    \end{scope}
  \end{tikzpicture}
  \caption{Top: Chadwick's experiment. Bottom: why a neutral particle with a
    proton's mass knocks out fast protons (head-on, elastic collisions).}\label{ch2:fig:chadwick}
\end{figure}

\textbf{Energy rules out gamma rays.} The fastest protons had about
\qty{5.7}{MeV} of kinetic energy ($\qty{1}{MeV} = \qty{1.602e-13}{J}$; see
\cref{ch1:sec:bigbang}). A gamma ray is a photon: it has energy but no mass. For
its energy, a photon carries very little momentum. So it is very bad at pushing
a heavy proton, like a ping-pong ball hitting a bowling ball. A photon would
need about \qty{50}{MeV} to give a proton \qty{5.7}{MeV}.

Could the source make such a photon? Suppose beryllium captured an alpha and gave
out a gamma ray
($\chTwoNuc{4}{2}{He} + \chTwoNuc{9}{4}{Be} \rightarrow \chTwoNuc{13}{6}{C} + \gamma$).
The photon could carry at most about \qty{14}{MeV}. A \qty{50}{MeV} photon from
this reaction would break conservation of energy. So the radiation was not
gamma rays.

\textbf{Momentum and kinetic energy give the mass.} Chadwick treated the
radiation as neutral particles of unknown mass $m$ and unknown speed $u$. They
make elastic collisions with the target nuclei.
\begin{itemize}
  \item For one head-on collision he wrote both laws: conservation of momentum
        and conservation of kinetic energy. Together they give the target's
        recoil speed in terms of $m$ and $u$.
  \item One target gives one equation, but there are two unknowns. So he used
        two targets. Hydrogen nuclei (about \qty{1}{u}) recoiled very fast.
        Nitrogen nuclei (about \qty{14}{u}) recoiled much more slowly.
  \item He divided one result by the other. The unknown speed $u$ cancelled.
        Only the mass was left: about \qty{1.15}{u}, close to a proton's mass.
\end{itemize}
(\qty{1}{u} $= \qty{1.661e-27}{kg}$, about one proton mass; see \cref{ch1:sec:fusion}.)
Snooker balls show the same idea: a moving ball passes the most energy to a ball
of equal mass (\cref{ch2:fig:chadwick}).

The reaction in the beryllium was
$\chTwoNuc{4}{2}{He} + \chTwoNuc{9}{4}{Be} \rightarrow \chTwoNuc{12}{6}{C} + \chTwoNuc{1}{0}{n}$.
Chadwick never observed a neutron directly. It has no charge, so it does not
ionise the gas in a detector and leaves no track. This is \term{indirect
evidence}, which means evidence from the effects of something that cannot be
detected directly. \Cref{ch2:tab:particles} sums up the atom.

\begin{table}
  \caption{The three particles of the atom (masses from the data sheet).}\label{ch2:tab:particles}
  \centering\small
  \begin{tabular}{@{}L{16mm}L{31mm}L{25mm}L{27mm}L{26mm}@{}}
    \toprule
    \colhead{Particle} & \colhead{Charge} & \colhead{Mass} & \colhead{Where} & \colhead{Found by}\\
    \midrule
    electron & $-e = -\qty{1.602e-19}{C}$ & \qty{9.109e-31}{kg} & around the nucleus & Thomson, 1897\\
    proton & $+e = +\qty{1.602e-19}{C}$ & \qty{1.673e-27}{kg} & in the nucleus & Rutherford, 1919\\
    neutron & 0 & \qty{1.675e-27}{kg} & in the nucleus & Chadwick, 1932\\
    \bottomrule
  \end{tabular}
\end{table}

\begin{formulacard}{Conservation of momentum and kinetic energy}
  \cardformula{\sum m\vec{v}_\mathrm{before} = \sum m\vec{v}_\mathrm{after}
    \qquad \sum\frac{1}{2}mv^2_\mathrm{before} = \sum\frac{1}{2}mv^2_\mathrm{after}}
  \begin{symbols}
    \sym{$m$}{mass of each object}{\si{kg}}
    \sym{$\vec{v}$}{velocity of each object, with a sign for its direction}{\si{m.s^{-1}}}
    \sym{$v$}{speed of each object (for kinetic energy)}{\si{m.s^{-1}}}
  \end{symbols}
  \inwords{in a collision the total momentum stays the same; in an elastic
    collision the total kinetic energy also stays the same.}
  \whentouse{momentum: any collision with no outside force. Kinetic energy: only
    elastic collisions, such as neutrons bouncing off nuclei.}
  \whennottouse{the kinetic-energy equation when objects stick together, or when
    energy turns into heat or sound.}
  \rearranged{object 1 hits object 2 head-on, and object 2 starts at rest
    ($u$ = velocity before, $v$ = velocity after): $m_1u_1 = m_1v_1 + m_2v_2$, so
    $v_1 = \dfrac{m_1u_1 - m_2v_2}{m_1}$}
  \tcbbreak % layout QA (book build): page break here, so "(continued)" is not a lone Trap line
  \tensecond{a neutron (\qty{1.675e-27}{kg}) at \qty{2.6e6}{m.s^{-1}} hits a
    nucleus at rest and bounces straight back at \qty{2.2e6}{m.s^{-1}}, so
    $v = -\qty{2.2e6}{m.s^{-1}}$. Momentum given to the nucleus
    $= m(u - v) = \num{1.675e-27}\times(\num{2.6e6} + \num{2.2e6})
    = \qty{8.0e-21}{kg.m.s^{-1}}$.}
  \trap{choose a positive direction first: a particle that bounces back has a
    negative velocity. Kinetic energy has no direction and is never negative.}
\end{formulacard}

\begin{workedexample}{Completing Chadwick's reaction}\label{ch2:we:equation}
Find the missing particle X: \
$\chTwoNuc{4}{2}{He} + \chTwoNuc{9}{4}{Be} \rightarrow \chTwoNuc{12}{6}{C} + \chTwoNuc{A}{Z}{X}$
\given{the equation, with three of the four nuclei known}
\find{$A$ and $Z$ for X, and what X is}
\formula{the totals of mass number and of atomic number are the same on both sides}
\substitute{$4 + 9 = 12 + A$, so $A = 1$; \quad $2 + 4 = 6 + Z$, so $Z = 0$}
\answer{X is \final{$\chTwoNuc{1}{0}{n}$}, a neutron.}
\checkstep{mass number 1 and no charge: a neutral particle with about a proton's
  mass, as Chadwick found \cmark.}
\end{workedexample}

\begin{workedexample}{A neutron hits a proton}\label{ch2:we:collision}
A neutron moving at \qty{2.0e7}{m.s^{-1}} hits a proton at rest head-on and stops.
Take both masses as \qty{1.67e-27}{kg} (the data-sheet values differ by only
0.1\%). Find the proton's speed, and show that kinetic energy is conserved.
\given{$m_\mathrm{n} = m_\mathrm{p} = \qty{1.67e-27}{kg}$, \
  $u_\mathrm{n} = \qty{2.0e7}{m.s^{-1}}$, \ $v_\mathrm{n} = 0$, \ $u_\mathrm{p} = 0$}
\find{the proton's speed $v_\mathrm{p}$; the kinetic energy before and after}
\formula{$\sum m\vec{v}_\mathrm{before} = \sum m\vec{v}_\mathrm{after}$, \quad $K = \tfrac12 mv^2$}
\substitute{Momentum, writing every term:
  $m_\mathrm{n}u_\mathrm{n} + m_\mathrm{p}(0) = m_\mathrm{n}(0) + m_\mathrm{p}v_\mathrm{p}$.
  The zero terms vanish, so $m_\mathrm{n}u_\mathrm{n} = m_\mathrm{p}v_\mathrm{p}$.
  Divide both sides by $m_\mathrm{p}$; the equal masses cancel:\par
  \eqlines{v_\mathrm{p} &= \frac{\qty{1.67e-27}{kg}\times\qty{2.0e7}{m.s^{-1}}}{\qty{1.67e-27}{kg}}
    = \qty{2.0e7}{m.s^{-1}}}\par
  Kinetic energy before (the neutron) and after (the proton):\par
  \eqlines{K_\mathrm{before} &= \tfrac12\times\qty{1.67e-27}{kg}\times(\qty{2.0e7}{m.s^{-1}})^2
    = \qty{3.34e-13}{J}\\
  K_\mathrm{after} &= \tfrac12\times\qty{1.67e-27}{kg}\times(\qty{2.0e7}{m.s^{-1}})^2
    = \qty{3.34e-13}{J}}}
\answer{\final{$v_\mathrm{p} = \qty{2.0e7}{m.s^{-1}}$}; $K = \qty{3.3e-13}{J}$
  before and after, so the collision is elastic.}
\checkstep{equal masses: the neutron hands all its speed to the proton, like
  snooker balls. This explains the fast protons from the wax. (With the exact
  data-sheet masses the neutron does not quite stop: it keeps about 0.06\% of
  its speed.)}
\end{workedexample}

\begin{workedexample}{A neutron bounces off carbon}\label{ch2:we:carbon}
A neutron (\qty{1.675e-27}{kg}) moving at \qty{2.50e7}{m.s^{-1}} hits a
carbon-12 nucleus (\qty{1.993e-26}{kg}) at rest, head-on. The carbon recoils at
\qty{3.876e6}{m.s^{-1}}. Find (a)~the neutron's velocity after the collision;
(b)~the fraction of the neutron's kinetic energy given to the carbon.
\given{right is positive; $m_\mathrm{n} = \qty{1.675e-27}{kg}$, \
  $m_\mathrm{C} = \qty{1.993e-26}{kg}$, \ $u_\mathrm{n} = +\qty{2.50e7}{m.s^{-1}}$, \
  $u_\mathrm{C} = 0$, \ $v_\mathrm{C} = +\qty{3.876e6}{m.s^{-1}}$}
\find{the neutron's velocity $v_\mathrm{n}$; the fraction $\dfrac{K_\mathrm{C}}{K_\mathrm{before}}$}
\formula{$\sum m\vec{v}_\mathrm{before} = \sum m\vec{v}_\mathrm{after}$, \quad $K = \tfrac12 mv^2$}
\substitute{Momentum before $=$ momentum after:
  $m_\mathrm{n}u_\mathrm{n} + 0 = m_\mathrm{n}v_\mathrm{n} + m_\mathrm{C}v_\mathrm{C}$.
  Subtract $m_\mathrm{C}v_\mathrm{C}$ from both sides, then divide by $m_\mathrm{n}$.\par
  \eqlines{v_\mathrm{n} &= \frac{m_\mathrm{n}u_\mathrm{n} - m_\mathrm{C}v_\mathrm{C}}{m_\mathrm{n}}\\
  m_\mathrm{n}u_\mathrm{n} &= \qty{1.675e-27}{kg}\times\qty{2.50e7}{m.s^{-1}}\\
  &= \qty{4.188e-20}{kg.m.s^{-1}}\\
  m_\mathrm{C}v_\mathrm{C} &= \qty{1.993e-26}{kg}\times\qty{3.876e6}{m.s^{-1}}\\
  &= \qty{7.725e-20}{kg.m.s^{-1}}\\
  v_\mathrm{n} &= \frac{\qty{4.188e-20}{kg.m.s^{-1}} - \qty{7.725e-20}{kg.m.s^{-1}}}{\qty{1.675e-27}{kg}}\\
  &= -\qty{2.112e7}{m.s^{-1}}}\par
  Kinetic energies:\par
  \eqlines{K_\mathrm{before} &= \tfrac12\times\qty{1.675e-27}{kg}\times(\qty{2.50e7}{m.s^{-1}})^2\\
  &= \qty{5.234e-13}{J}\\
  K_\mathrm{C} &= \tfrac12\times\qty{1.993e-26}{kg}\times(\qty{3.876e6}{m.s^{-1}})^2\\
  &= \qty{1.497e-13}{J}\\
  \text{fraction} &= \frac{\qty{1.497e-13}{J}}{\qty{5.234e-13}{J}} = 0.286}}
\answer{(a)~\final{$v_\mathrm{n} = \qty{2.11e7}{m.s^{-1}}$ backwards} (the neutron
  bounces back). \ (b)~\final{about 29\%}.}
\checkstep{the neutron's kinetic energy after is
  $\tfrac12\times\qty{1.675e-27}{kg}\times(\qty{2.112e7}{m.s^{-1}})^2 = \qty{3.736e-13}{J}$.
  Total after $= \qty{1.497e-13}{J} + \qty{3.736e-13}{J} = \qty{5.233e-13}{J}$,
  the same as before (to rounding), so the collision is elastic \cmark. A heavy target takes
  much less energy than a proton does.}
\end{workedexample}

\begin{mistakes}
\begin{itemize}
  \item \xmark\ ``Chadwick used momentum alone to find the mass.'' \cmark\ He
        needed momentum \emph{and} kinetic energy, with two targets.
  \item \xmark\ ``Fields did not bend it, so it was gamma rays.'' \cmark\ That only
        shows it has no charge.
  \item \xmark\ Writing the neutron as $\chTwoNuc{1}{1}{n}$. \cmark\ It is
        $\chTwoNuc{1}{0}{n}$ (charge 0).
\end{itemize}
\end{mistakes}

\begin{examtip}
``Assess how important the conservation laws were in Chadwick's identification
of the neutron'' (5 marks). Full marks need:
\begin{enumerate}
  \item both laws, stated;
  \item energy: a gamma ray would need about \qty{50}{MeV}, more than the
        reaction gives, so that idea was rejected;
  \item momentum and kinetic energy, with H and N targets, gave a mass close to
        the proton's;
  \item a judgement: ``The laws were essential, because the neutron could not be
        detected directly.''
\end{enumerate}
\end{examtip}

\begin{quickcheck}
  \item Complete and identify:
        $\chTwoNuc{4}{2}{He} + \chTwoNuc{14}{7}{N} \rightarrow \chTwoNuc{17}{8}{O} + {?}$
  \item The neutral radiation from beryllium could not be detected directly.
        Explain how a block of paraffin wax in its path let Chadwick measure its effects.
  \item A neutron (\qty{1.675e-27}{kg}, \qty{3.00e7}{m.s^{-1}}) hits a nitrogen
        nucleus (\qty{2.325e-26}{kg}) at rest head-on; the nitrogen recoils at
        \qty{4.03e6}{m.s^{-1}}. Find (a)~the neutron's velocity afterwards;
        (b)~the fraction of kinetic energy the nitrogen got.
  \item A classmate says the radiation from beryllium was gamma rays, because
        fields did not bend it. Use conservation of energy to explain why this
        idea fails.
\end{quickcheck}

%% ==================================================================
\section{Key experiments}\label{ch2:sec:experiments}
\covers{M8.IQ2.01, M8.IQ2.02}

Exam questions may ask you to \textbf{model} an experiment. This means drawing
and labelling its set-up, or explaining how a computer simulation of it works.

Questions also ask about three qualities. An experiment is \term{valid} when it
tests what it claims to test. It is \term{reliable} when repeating it gives the
same results. It is \term{accurate} when its result is close to the true value.
To answer, name one feature of the experiment. Then say what it controls or
improves.

\begin{experiment}{Early cathode-ray experiments (1870s--1890s)}
\aim{to find whether cathode rays are waves or particles.}
\apparatus{evacuated tubes (tubes with almost all the air pumped out),
  high-voltage DC supply, Maltese cross, paddle wheel, charged plates, magnet
  (\cref{ch2:fig:crt,ch2:fig:tubes} on pages~\pageref{ch2:fig:crt} and~\pageref{ch2:fig:tubes}).}
\method{\begin{enumerate}
  \item Put one object in the beam, or bring charged plates or a magnet near the tube.
  \item Watch the fluorescent end of the tube and record what happens.
\end{enumerate}
\textbf{Safety:} the high voltage can give a dangerous shock, so never touch the
terminals. The tube can give off X-rays, so run it only for short times. The
glass could implode (break inwards), so wear safety glasses.}
\results{\cref{ch2:tab:evidence}.}
\conclusion{the rays travel in straight lines and carry negative charge: they are
  negative particles.}
\significance{showed that the rays are negative particles, and set up Thomson's
  1897 measurement, which ended the 20-year debate.}
\validity{\begin{itemize}
  \item \textbf{Validity:} Hertz's poor vacuum made his plate test invalid;
        heating, not momentum, mostly turned the paddle wheel.
  \item \textbf{Reliability:} many labs, tubes and gases gave the same results.
  \item \textbf{Accuracy:} almost no numbers were measured (only a rough speed),
        so the tests showed properties, not values.
\end{itemize}}
\end{experiment}

\begin{experiment}{Thomson's charge-to-mass experiment (1897)}
\aim{to measure $q/m$ for cathode-ray particles.}
\apparatus{tube with cathode, collimating anodes, plates, coils and a
  fluorescent screen with a scale (\cref{ch2:fig:thomson} on page~\pageref{ch2:fig:thomson}).
  A school e/m tube (\cref{ch2:fig:emtube}) models the same idea.}
\method{\begin{enumerate}
  \item Switch on the high voltage so a narrow beam hits the middle of the screen.
  \item Switch on both fields. Adjust them until the beam goes straight. Then
        $qE = qvB$, so $v = \dfrac{E}{B}$.
  \item Switch off the electric field. Use how far the spot moves to find the
        radius $r$.
  \item Calculate $\dfrac{q}{m} = \dfrac{E}{B^2 r}$.
  \item Repeat with other cathode metals, gases and voltages.
\end{enumerate}}
\results{about \qty{e11}{C.kg^{-1}} (today \qty{1.759e11}{C.kg^{-1}}) for every metal and gas.}
\conclusion{identical negative particles, far lighter than atoms, in all matter:
  electrons.}
\significance{the first subatomic particle; it led to the plum pudding model.}
\validity{\begin{itemize}
  \item \textbf{Validity:} high vacuum; uniform fields at \ang{90} over the same
        region; narrow beam; gravity is negligible.
  \item \textbf{Reliability:} the same $q/m$ for many metals, gases and voltages.
  \item \textbf{Accuracy:} only rough, because $r$ came from the spot's movement
        and the fields are not uniform at their edges.
\end{itemize}}
\end{experiment}

\begin{experiment}{Millikan's oil drop experiment (1909--1913)}
\aim{to find whether charge comes in a smallest unit, and measure it.}
\apparatus{atomiser, horizontal plates, adjustable DC supply, microscope, light,
  X-ray source (\cref{ch2:fig:millikan} on page~\pageref{ch2:fig:millikan}).}
\method{\begin{enumerate}
  \item Spray oil drops into the chamber and watch one drop through the microscope.
  \item With the voltage off, time the drop falling at its terminal velocity.
        This speed gives its mass.
  \item Switch on the voltage and adjust $V$ until the drop is still. Calculate
        $q = \dfrac{mgd}{V}$.
  \item Change the drop's charge with X-rays and repeat.
  \item Repeat for many drops.
\end{enumerate}}
\results{every charge was a whole number times \qty{1.592e-19}{C} (today
  \qty{1.602e-19}{C}).}
\conclusion{charge is quantised; the smallest unit $e$ is the electron's charge.}
\significance{it gave $e$, and with Thomson's $q/m$ the electron's mass.}
\validity{\begin{itemize}
  \item \textbf{Validity:} closed chamber (no air currents); oil does not evaporate.
  \item \textbf{Reliability:} many drops, and the same drop with different charges.
  \item \textbf{Accuracy:} his value was 0.6\% too low. He used a slightly wrong
        value for the viscosity of air (how strongly air resists things moving
        through it). This is a systematic error: an error that shifts every
        result the same way.
\end{itemize}}
\end{experiment}

\begin{experiment}{The Geiger--Marsden experiment (1909)}
\aim{to test the plum pudding model by firing alpha particles at thin metal foil.}
\apparatus{alpha source in a lead block, thin gold foil, zinc sulfide screen and
  microscope that turn around the foil, vacuum chamber
  (\cref{ch2:fig:goldfoil} on page~\pageref{ch2:fig:goldfoil}).}
\method{\begin{enumerate}
  \item Fire a narrow beam of alphas at the foil in a vacuum chamber.
  \item In the dark, count the flashes on the screen through the microscope.
  \item Move the microscope to a new angle and count again.
\end{enumerate}}
\results{most went straight through; some small deflections; about 1 in 8000
  turned by more than \ang{90}.}
\conclusion{a tiny, dense, positive nucleus; the atom is mostly empty space.}
\significance{the discovery of the nucleus, the basis of Bohr's model and of
  nuclear physics.}
\validity{\begin{itemize}
  \item \textbf{Validity:} the chamber is a vacuum, so air does not scatter the
        alphas. The foil is very thin, so any large deflection comes from one
        close pass by a single nucleus. Gold has a large nuclear charge ($+79e$)
        and can be made very thin.
  \item \textbf{Reliability:} many thousands of flashes, counted by more than one observer.
  \item \textbf{Accuracy:} faint flashes are easy to miss. Very few alphas turn by
        large angles, so those counts have the largest uncertainty.
\end{itemize}}
\end{experiment}

\begin{experiment}{Chadwick's discovery of the neutron (1932)}
\aim{to identify the neutral radiation from beryllium hit by alpha particles.}
\apparatus{polonium source, beryllium, paraffin wax and other targets, ionisation
  chamber (\cref{ch2:fig:chadwick} on page~\pageref{ch2:fig:chadwick}).}
\method{\begin{enumerate}
  \item Let alphas from polonium hit beryllium.
  \item Let the radiation hit paraffin wax. Measure the energy of the protons
        knocked out.
  \item Repeat with other targets, such as nitrogen.
  \item Apply conservation of energy and momentum.
\end{enumerate}}
\results{protons up to \qty{5.7}{MeV}; nitrogen nuclei recoil (move off after
  being hit) slowly; no single gamma-ray energy fits both.}
\conclusion{neutral particles with about the proton's mass: neutrons.}
\significance{completed our picture of the atom and explained the missing mass.
  Neutrons were later used to split nuclei (\cref{ch4:sec:fission}).}
\validity{\begin{itemize}
  \item \textbf{Validity:} neutrons were found only indirectly, so the result
        depends on using the conservation laws correctly.
  \item \textbf{Reliability:} recoils from several different target nuclei all
        fitted one neutron mass. A second method (the energy balance of a boron
        reaction) gave a similar mass, about \qty{1.007}{u}.
  \item \textbf{Accuracy:} the first estimate (about \qty{1.15}{u}) had about 10\%
        uncertainty.
\end{itemize}}
\end{experiment}
